% ================================================================
%  第 6 章  Plateau 问题的嵌入解
% ================================================================
\chapter{Plateau 问题的嵌入解}
\label{chap:6}

{\setlength{\emergencystretch}{2em}%
在一般性地讨论\newterm{unique-continuation}{唯一延拓}{unique continuation}与零点集之后, 我们研究极小曲面在一个分支点附近或
一个非横截相交点附近的局部描述. 遵循 Osserman 与 Gulliver 的做法, 我们
排除 Plateau 问题解的内部分支点. 在本章余下的部分, 我们证明: 当边界位于
某个\newterm{mean-convex}{平均凸}{mean convex}区域的边界上时, Plateau 问题的解是嵌入的. 这最后一个结果属于
Meeks 与 Yau.\par}

% ================================================================
\section{唯一延拓}
\label{sec:6.1}

在考察极小曲面的局部性质时, 对其可能出现的自交与分支点有一个好的描述
将会很有用. 这可以借助关于椭圆方程解的强唯一延拓的结论来实现. 后面两节
的结果将在其后的各节中用到: 先用它们证明极小圆盘是浸入的, 然后 (在附加
假设之下) 证明它们是嵌入的.

具体地说, 我们需要下述定量形式的强唯一延拓:

\begin{theorem}[定量强唯一延拓]
\label{thm:6.1}
设 $v$ 是 $D$ 上的非常数解, 满足
\begin{equation}
0=(a_{i,j}v_{x_j})_{x_i}+b_iv_{x_i}+cv,
\label{eq:6.1}
\end{equation}
其中 $a_{ij}$ 对称、一致椭圆且 Lipschitz, $b_i,c$ 连续. 则存在
$\epsilon=\epsilon(a_{ij},b_i,c)>0$ 与 $\bar{d}$, 使得对任意的
$2t<\epsilon$,
\begin{equation}
\frac{\int_{|x|=2t}v^2}{\int_{|x|=t}v^2}\leq 2^{2\bar{d}+1}.
\label{eq:6.2}
\end{equation}
特别地, 存在某个整数 $d\leq\bar{d}$, 使得 $v$ 在原点处恰消没 $d$ 阶.
\end{theorem}

事实上, 由定理~\ref{thm:6.1} 的证明将看出, $\bar{d}$ 可由 $v$
自 $D_\epsilon$ 到 $D_{1/2}$ 的增长给出上界估计.

\begin{remark}
\label{rem:6.2}
特别地, 定理~\ref{thm:6.1} 蕴含 $v$ 不可能在任何点处以无穷阶消没. 具有
这一性质的算子称为满足\emph{强唯一延拓}. 注意: 若 $L$ 具有这一强唯一
延拓性质, 且 $Lv=0$ 于连通区域 $\Omega$ 上, 而 $v$ 在 $\Omega$ 的某个
开子集上为零, 则必有 $v\equiv0$ 于 $\Omega$ 上.
\end{remark}

在 $a_{ij}=\delta_{ij}$ 且 $b_i=c=0$ 的情形, 该定理与复分析中的
Hadamard 三圆定理密切相关.

\mtpara{Almgren 的频率函数}
我们先利用 Almgren 的频率函数, 在 $\R^n$ 上 Laplacian 的情形证明定理~\ref{thm:6.1}; 这个证明
很容易修改以适用于一般情形.

设 $\Delta u=0$ 于 $\R^n$ 上, 并记 $u_r=\frac{\partial u}{\partial r}$ 为其
径向导数. 给定 $r>0$, 定义尺度不变的量 $I(r)$、$E(r)$ 以及\newterm{frequency}{频率函数}{frequency function}
$U(r)$ (参看 Almgren, \cite{Am2}):
\begin{align}
I(r)&=r^{1-n}\int_{\partial B_r}u^2,\label{eq:6.3}\\
E(r)&=r^{2-n}\int_{B_r}|\nabla u|^2=r^{2-n}\int_{\partial B_r}uu_r,
  \label{eq:6.4}\\
U(r)&=\frac{E(r)}{I(r)}.\label{eq:6.5}
\end{align}

对 $I$ 求导给出
\begin{equation}
I'(r)=2r^{1-n}\int_{\partial B_r}uu_r=\frac{2E(r)}{r},
\label{eq:6.6}
\end{equation}
其中最后一步用了散度定理, 因为对调和函数 $u$ 有
$\dive(u\nabla u)=|\nabla u|^2$.

\begin{lemma}
\label{lem:6.3}
$E$ 的导数为
\begin{equation}
E'(r)=2r^{2-n}\int_{\partial B_r}|u_r|^2.
\label{eq:6.7}
\end{equation}
\end{lemma}

当 $u$ 是二维圆盘上的调和函数时, 引理~\ref{lem:6.3} 可以通过写
$u=\operatorname{Re}f$ (其中 $f$ 全纯)、利用 Cauchy--Riemann 方程并
分部积分来证明.

\begin{proof}[引理~\ref{lem:6.3} 的证明]
利用余面积公式 (即 \eqref{eq:1.59}) 对 $E$ 求导, 给出
\begin{equation}
E'(r)=\frac{2-n}{r}E(r)+r^{2-n}\int_{\partial B_r}|\nabla u|^2.
\label{eq:6.8}
\end{equation}
在 $B_r\subset\R^n$ 上置 $r^2=|x|^2$, 于是 Hessian 为
$(r^2)_{ij}=2\delta_{ij}$; 注意
\begin{equation}
\dive\Bigl[(r^2)_ju_ju_k-\frac12u_j^2(r^2)_k\Bigr]
=(2-n)|\nabla u|^2+2ru_r\Delta u.
\label{eq:6.9}
\end{equation}
对 \eqref{eq:6.9} 用散度定理得
\begin{equation}
(2-n)\int_{B_R}|\nabla u|^2=2R\int_{\partial B_R}u_r^2
-R\int_{\partial B_R}|\nabla u|^2.
\label{eq:6.10}
\end{equation}
把 \eqref{eq:6.10} 代入 $E'(r)$ 的表达式即得引理.
\end{proof}

\begin{lemma}
\label{lem:6.4}
对 $\log U(r)$ 求导给出
\begin{equation}
\frac{d}{dr}\log U(r)=\frac{2r^{-n}}{E(r)}
\int_{\partial B_r}\bigl(ru_r-U(r)u\bigr)^2\geq0.
\label{eq:6.11}
\end{equation}
\end{lemma}

\begin{proof}
对 $\log U(r)$ 求导, 再代入 $I$ 与 $E$ 的导数\footnote{原书此处作
``the derivatives of $I$ and $D$''; 本章只出现 $I(r)$、$E(r)$、$U(r)$
三个量, 故 $D$ 应为 $E$ 之误.}, 给出
\begin{align}
\frac{d}{dr}\log U(r)
&=\frac{E'(r)}{E(r)}-\frac{I'(r)}{I(r)}\notag\\
&=\frac{2}{E(r)}\Bigl(r^{2-n}\int_{\partial B_r}|u_r|^2
-\frac{E^2(r)}{rI(r)}\Bigr).
\label{eq:6.12}
\end{align}
注意
\begin{equation}
\int_{\partial B_r}\bigl(ru_r-U(r)u\bigr)^2
=\int_{\partial B_r}r^2|u_r|^2-\frac{rE^2(r)}{I(r)}.
\label{eq:6.13}
\end{equation}
把它代入 \eqref{eq:6.12} 即得引理.
\end{proof}

\begin{example}
\label{ex:6.5}
设 $v(x)$ 是 $d$ 次齐次调和多项式. 容易看出 $rv_r=dv$, 从而
$E(t)=dI(t)$. 于是此时 $U(t)\equiv d$. 更一般地, 设
$w=w_1+w_2$ 是调和函数, 其中 $w_1$ 是 $d$ 次齐次的, 而 $w_2$ 的
Taylor 级数从 $d+1$ 次开始. 利用球面调和函数的正交性可知, 此时
$U(t)\geq d$.
\end{example}

尽管不会用到, 我们仍指出: 若 $v$ 调和且频率为常数, 则 $v$ 必是齐次的,
从而是球面调和函数. 确切地说, 有下述结论:

\begin{lemma}
\label{lem:6.6}
若 $v$ 在 $D$ 上调和且频率函数 $U(t)$ 为常数 (比如 $U(t)=d$), 则 $v$
是 $d$ 次齐次调和多项式.
\end{lemma}

\begin{proof}
由于 $U(t)=d$ 为常数, 引理~\ref{lem:6.4} 给出
\begin{equation}
v_r(x)=d|x|^{-1}v(x).
\label{eq:6.14}
\end{equation}
由 \eqref{eq:6.14} 即知 $v$ 是 $d$ 次齐次的, 从而是球面调和函数.
\end{proof}

现在可以证明 $\Delta v=0$ 情形的定理~\ref{thm:6.1} 了. 事实上, 此时我们将
证明一个更强的结论, 它将涵盖下文所述的定理~\ref{thm:6.12}.

\begin{proof}[$\Delta v=0$ 时定理~\ref{thm:6.1} 的证明]
我们将由频率函数的单调性来证明此结论. 对 $\log I(t)$ 求导 (见
\eqref{eq:6.6}) 给出
\begin{equation}
(\log I)'(t)=2\frac{U(t)}{t}.
\label{eq:6.15}
\end{equation}
于是利用 $U(t)$ 的单调性可写
\begin{equation}
\frac{I(1)}{I(\frac12)}
=\mathrm{e}^{2\int_{\frac12}^{1}\frac{U(t)}{t}dt}
\geq\mathrm{e}^{2U(\frac12)\log2}=4^{U(\frac12)}.
\label{eq:6.16}
\end{equation}
再次利用 $U$ 的单调性, 由 \eqref{eq:6.16} 看出 $U(t)$ 在
$t\leq\frac12$ 上一致有界. 因此, 对 \eqref{eq:6.15} 积分并利用 $U(t)$
的一致界, 当 $0<2s\leq t<\frac14$ 时,
\begin{equation}
\frac{I(2s)}{I(s)}
=\mathrm{e}^{2\int_s^{2s}\frac{U(t')}{t'}dt'}
\leq2^{2U(2s)}\leq\frac{I(2t)}{I(t)}.
\label{eq:6.17}
\end{equation}
估计 \eqref{eq:6.17} 称为\newterm{doubling}{加倍估计}{doubling estimate}. 迭代 \eqref{eq:6.17} 便给出
原点处消没阶数的界 (这就给出定量强唯一延拓). 因此, $v$ 在原点处的
Taylor 多项式不恒为零; 记 $v^1$ 为第一个非零 Taylor 多项式 (其次数为
$d\geq1$). 由于 $v$ 调和, $v_1$ 与 $v-v_1$ 均调和. 按构造, $v-v_1$
至少消没 $d+1$ 阶, 定理得证.
\end{proof}

\mtpara{一般算子的频率函数}
证明这一强唯一延拓定理的关键在于得到加倍估计 \eqref{eq:6.17}. 该加倍
估计是频率的一致界的直接推论. 当 $\Delta v=0$ 于圆盘上时, 我们由频率的
单调性推出了这一界. 对更一般的方程, 如定理~\ref{thm:6.1} 中的方程, 我们
同样可以写出一个频率函数 $U(t)$, 但它未必单调. 不过, 只要系数足够正则,
就存在某个常数 $\Lambda$ 使 $\mathrm{e}^{\Lambda t}U(t)$ 单调. 这已足以
在 $t$ 较小时建立 $U(t)$ 的一致界, 从而得到加倍估计.

作为例子, 考虑 $a_{ij}=f\delta_{ij}$ 的情形, 其中 $f(0)=1$、
$|\nabla f|\leq\lambda$、$b_i=c=0$. 设 $v$ 是 $D$ 上的非常数解, 满足
\begin{equation}
\dive(f\nabla v)=0.
\label{eq:6.18}
\end{equation}
当 $b_i$ 与 $c$ 不为零时如何修改这些论证, 对读者是清楚的. 此外, 形如
\eqref{eq:6.18} 的方程 (包括 $b_i$ 与 $c$ 出现的情形) 已足以覆盖我们
的全部应用.

\begin{remark}
\label{rem:6.7}
如 \cite{GaLf} 第 2 节末所指出的, 形如 \eqref{eq:6.1} 的方程可以通过
作适当的度量变换化为形如 \eqref{eq:6.18} 的方程 (细节见 \cite{GaLf}).
\end{remark}

一般说来, 上述变换需要引入新度量, 从而带来扰动项. 在二维情形, 这可利用
等温坐标的存在性来避免.

\begin{example}[Riemann 度量下的 Laplacian]
\label{ex:6.8}
若 $g_{ij}$ 是 $\R^2$ 某坐标卡中的 Riemann 度量, 则相应的 Laplacian
$\Delta_g$ 为
\begin{equation}
\Delta_gv=\frac{1}{\sqrt{g}}\bigl(\sqrt{g}g^{ij}v_{x_j}\bigr)_{x_i},
\label{eq:6.19}
\end{equation}
其中 $g$ 与 $g^{ij}$ 分别是 $g_{ij}$ 的行列式与逆矩阵. 等温坐标的存在性
使我们可以作坐标变换后设 $g_{ij}=f\delta_{ij}$, 其中 $f$ 是正函数. 此时
\eqref{eq:6.19} 变为
\begin{equation}
\Delta_gv=\frac1f v_{x_ix_i}.
\label{eq:6.20}
\end{equation}
对 $0<t\leq1$, 定义 $I(t)$ 与 $E(t)$ 为
\begin{equation}
I(t)=t^{-1}\int_{|x|=t}fv^2.
\label{eq:6.21}
\end{equation}
利用 Stokes 定理与公式 $\dive(f\nabla v^2)=2f|\nabla v|^2$ (由
\eqref{eq:6.18} 得到), 有
\begin{equation}
E(t)=\int_{|x|\leq t}f|\nabla v|^2=\int_{|x|=t}fv_rv.
\label{eq:6.22}
\end{equation}
\end{example}

我们再一次按 \eqref{eq:6.5} 的方式定义频率 $U=I/E$.

此后我们总假设 $t$、$|x|\leq\frac{1}{2\lambda}$, 于是
$\frac12\leq f(x)\leq\frac32$. 对 $I(t)$ 求导给出
\begin{equation}
I'(t)=2t^{-1}\int_{|x|=t}fvv_r+t^{-1}\int_{|x|=t}f_rv^2,
\label{eq:6.23}
\end{equation}
于是 Stokes 定理与界 $|\nabla f|\leq2\lambda f$ 给出
\begin{equation}
\Bigl|I'(t)-2\frac{E(t)}{t}\Bigr|\leq2\lambda I(t).
\label{eq:6.24}
\end{equation}
由 $U(t)$ 的定义与 \eqref{eq:6.24} 得到微分不等式
\begin{equation}
\Bigl|(\log I)'(t)-2\frac{U(t)}{t}\Bigr|\leq2\lambda.
\label{eq:6.25}
\end{equation}
余面积公式 (即 \eqref{eq:1.59}) 给出
\begin{equation}
E'(t)=\int_{|x|=t}f|\nabla v|^2.
\label{eq:6.26}
\end{equation}
像前面那样对 \eqref{eq:6.26} 应用 Stokes 定理, 得到
\begin{align}
E'(t)&=\int_{|x|=t}f|\nabla v|^2
=t^{-1}\int_{|x|=t}\bigl\langle f|\nabla v|^2x,|x|^{-1}x\bigr\rangle
\notag\\
&=2t^{-1}E(t)+t^{-1}\int_{|x|\leq t}x_j(fv_{x_i}v_{x_i})_{x_j},
\label{eq:6.27}
\end{align}
这是因为
\begin{equation}
\dive(f|\nabla v|^2x)=(fv_{x_i}^2x_j)_{x_j}
=x_j(fv_{x_i}v_{x_i})_{x_j}+2fv_{x_i}^2.
\label{eq:6.28}
\end{equation}
由于 $(fv_{x_i})_{x_i}=0$、$v_{x_ix_j}=v_{x_jx_i}$ 且
$(x_i)_{x_j}=\delta_{ij}$,
\begin{align}
x_j(fv_{x_i}v_{x_i})_{x_j}
&=x_jf_{x_j}v_{x_i}v_{x_i}+2x_jfv_{x_i}v_{x_jx_i}\notag\\
&=x_jf_{x_j}v_{x_i}v_{x_i}+2(x_jfv_{x_i}v_{x_j})_{x_i}
-2x_jv_{x_j}(fv_{x_i})_{x_i}-2fv_{x_i}^2\notag\\
&=x_jf_{x_j}v_{x_i}v_{x_i}+2(x_jfv_{x_i}v_{x_j})_{x_i}-2fv_{x_i}^2.
\label{eq:6.29}
\end{align}
把式 \eqref{eq:6.29} 代入 \eqref{eq:6.27}, 分部积分, 并利用界
$|\nabla f|\leq2\lambda f$, 得到
\begin{equation}
\Bigl|E'(t)-2\int_{|x|=t}fv_r^2\Bigr|\leq4\lambda E(t).
\label{eq:6.30}
\end{equation}
结合 \eqref{eq:6.25} 与 \eqref{eq:6.30}, 对 $\log U(t)$ 求导给出
\begin{equation}
\Bigl|(\log U)'(t)-\frac{2}{E(t)}\int_{|x|=t}fv_r^2
+2\frac{U(t)}{t}\Bigr|\leq6\lambda.
\label{eq:6.31}
\end{equation}
利用 Cauchy--Schwarz 不等式给出
\begin{equation}
0\leq\frac{2}{E(t)}\int_{|x|=t}fv_r^2-2\frac{U(t)}{t}.
\label{eq:6.32}
\end{equation}
结合 \eqref{eq:6.31} 与 \eqref{eq:6.32} 可知, $\mathrm{e}^{6\lambda t}U(t)$
单调不减. 特别地, 对一切 $0<t\leq\frac{1}{2\lambda}$,
\begin{equation}
U(t)\leq\mathrm{e}^{6\lambda t}U(t)\leq\mathrm{e}^{3}
U\Bigl(\frac{1}{2\lambda}\Bigr).
\label{eq:6.33}
\end{equation}
把频率的一致界 \eqref{eq:6.33} 与微分不等式 \eqref{eq:6.25} 结合起来,
即得 $v$ 的消没阶数所需的界. 至此, 对形如 \eqref{eq:6.18} 的方程,
定理~\ref{thm:6.1} 的证明完成.

% ================================================================
\section{零点集与临界集的局部描述}
\label{sec:6.2}

利用前一节的结果, 我们将得到圆盘 $D$ 上椭圆方程解的一个有用的局部描述.
我们的重点主要在于得到局部渐近展开, 以及这些展开对零点集与临界集的
推论.

\begin{definition}[零点集与临界集]
\label{def:6.9}
给定函数 $u$, $u$ 的\newterm{nodal-set}{零点集}{nodal set}是使 $u$ 为零的集合 $\{u=0\}$. $u$ 的
\newterm{critical-set}{临界集}{critical set}是使 $u$ 与 $\nabla u$ 同时为零的集合
$\{u=0\}\cap\{\nabla u=0\}$.
\end{definition}

在进入本节的主要结果之前, 我们需要一些关于椭圆估计的背景材料. 正如在
第 1 章中所见, 平均值不等式用调和函数在较大球上的平均值来控制其最大值.
把它与反向 Poincar\'e 不等式结合起来, 便可同样地控制所有高阶导数. 对更
一般的方程, 类似的论证给出下述引理 (证明见 \cite{GiTr} 的定理 8.8 或
\cite{HaLi}):

\begin{lemma}
\label{lem:6.10}
设 $v$ 是 \eqref{eq:6.1} 在 $D$ 上的解, 并设 $a_{ij},b_i,c$ 光滑. 则存在
仅依赖于 $a_{ij},b_i,c$ 的常数 $\epsilon>0$ 与 $C<\infty$, 使得当
$|x|\leq\epsilon$ 时,
\begin{equation}
|v(x)|^2+|x|^2|\nabla v(x)|^2+|x|^4|\nabla^2v(x)|^2
\leq C|x|^{-2}\int_{B_{2|x|}}v^2.
\label{eq:6.34}
\end{equation}
\end{lemma}

把定理~\ref{thm:6.1} 与引理~\ref{lem:6.10} 结合起来, 得到下述推论:

\begin{corollary}
\label{cor:6.11}
设 $v$ 是 \eqref{eq:6.1} 在 $D$ 上的解, $\bar{d}$ 是定理~\ref{thm:6.1}
给出的消没阶数之界. 再设 $a_{ij}$、$b_i$ 与 $c$ 光滑. 则存在仅依赖于
$a_{ij}$、$b_i$、$c$ 与 $\bar{d}$ 的常数 $\epsilon>0$ 与 $C<\infty$,
使得当 $|x|\leq\epsilon$ 时,
\begin{equation}
|v(x)|^2+|x|^2|\nabla v(x)|^2+|x|^4|\nabla^2v(x)|^2\leq CI(|x|).
\label{eq:6.35}
\end{equation}
\end{corollary}

\begin{proof}
由定理~\ref{thm:6.1}, 频率 $U(t)$ 在某个球 $\{|x|\leq2\epsilon\}$ 上被
$\bar{d}$ 一致界定. 因此, 对任意的 $s\leq\epsilon$ 有
\begin{equation}
I(2s)\leq2^{2\bar{d}}I(s).
\label{eq:6.36}
\end{equation}
对界 \eqref{eq:6.36} 积分 (并利用 $I(t)$ 的平凡单调性) 给出
\begin{equation}
\int_{|x|\leq2s}v^2=\int_{t=0}^{2s}\int_{|x|=t}v^2dt
\leq\int_{t=0}^{2s}tI(t)dt\leq C_1s^2I(s).
\label{eq:6.37}
\end{equation}
把 \eqref{eq:6.34} 与 \eqref{eq:6.37} 结合起来即得估计 \eqref{eq:6.35}.
\end{proof}

\begin{theorem}[局部渐近展开, \cite{Ber}]
\label{thm:6.12}
设 $v$ 是 \eqref{eq:6.1} 的解, $\bar{d}$ 与 $d$ 分别是定理~\ref{thm:6.1}
给出的消没阶数之界与在 $0$ 处的消没阶数. 再设 $a_{ij},b_i$ 与 $c$
光滑. 则存在线性变换 $T$, 使得在经 $T$ 旋转之后可写
\begin{equation}
v(x)=p_d(x)+q(x),
\label{eq:6.38}
\end{equation}
其中 $p_d$ 是 $d$ 次齐次调和多项式, 且
\begin{equation}
|q(x)|+|x||\nabla q(x)|+\cdots+|x|^d|\nabla^dq(x)|\leq C|x|^{d+1}
\label{eq:6.39}
\end{equation}
对某个仅依赖于 $a_{ij},b_i,c$ 与 $\bar{d}$ 的 $C<\infty$ 成立.
\end{theorem}

\begin{proof}
经一个线性变换, 可设 $a_{ij}(0)=\delta_{ij}$, 且
$a_{ij}(x)-\delta_{ij}=c_{ij}(x)$, 其中 $|c_{ij}(x)|\leq\lambda|x|$. 于是
可把 \eqref{eq:6.1} 写成
\begin{align}
\Delta v&=-(c_{i,j}v_{x_j})_{x_i}-b_iv_{x_i}-cv\notag\\
&=-c_{i,j}v_{x_ix_j}-b_i'v_{x_i}-cv,
\label{eq:6.40}
\end{align}
其中 $b_i'=b_i+(c_{i,j})_{x_j}$.

由定理~\ref{thm:6.1}, $v(x)$ 在原点处恰消没 $d$ 阶. 因此极限
\begin{equation}
\lim_{t\to0}t^{-2d}I(t)
\label{eq:6.41}
\end{equation}
存在且不为零. 事实上, 若 $p_d(x)$ 表示 $v$ 在原点处的 $d$ 次 Taylor
多项式, 则
\begin{equation}
\lim_{t\to0}t^{-2d}I(t)=\lim_{t\to0}t^{-2d-1}\int_{|x|=t}p_d^2.
\label{eq:6.42}
\end{equation}
定义误差项 $q$ 为 $q(x)=v(x)-p_d(x)$, 于是 $q$ 在原点处至少消没 $d+1$ 阶.

由于 $v$ 满足 \eqref{eq:6.40}, 有
\begin{equation}
\lim_{|x|\to0}|x|^{d-2}\Delta v
=-\lim_{|x|\to0}|x|^{d-2}\bigl(c_{i,j}v_{x_ix_j}-b_i'v_{x_i}-cv\bigr).
\label{eq:6.43}
\end{equation}
频率的界蕴涵估计 \eqref{eq:6.35} 适用于 $v$. 因此 \eqref{eq:6.43} 中后
两项的极限均为零, 且 $v_{x_ix_j}$ 是 $|x|^{d-2}$ 阶的. 然而
$|c_{ij}|\leq\lambda|x|$, 故右端最后一项也趋于零. 于是
\begin{equation}
|\Delta v|\leq C|x|^{d-1}.
\label{eq:6.44}
\end{equation}
另一方面, 由于 $v-p_d$ 至少是 $d+1$ 阶的, Hessian $\nabla^2(v-p_d)$
至少是 $d-1$ 阶的. 因此 \eqref{eq:6.44} 蕴涵 $\Delta p_d$ 至少是
$d-1$ 阶的.

然而 $p_d$ 是 $d$ 次齐次的, 从而 $\Delta p_d$ 是 $d-2$ 次齐次多项式.
一个至少为 $d-1$ 阶的 $d-2$ 次齐次多项式必须恒为零, 故 $p_d$ 是球面
调和函数.

把 $q$ 的定义代入 \eqref{eq:6.40}, 看出 $q$ 本身满足一个椭圆方程
\begin{equation}
\Delta q=\Delta(v-p_d)=\Delta v
=-c_{i,j}v_{x_ix_j}-b_i'v_{x_i}-cv.
\label{eq:6.45}
\end{equation}
如前所验证, \eqref{eq:6.45} 的右端至少消没 $d+1$ 阶. 于是由
引理~\ref{lem:6.10} 即得 \eqref{eq:6.39} 中的高阶估计.
\end{proof}

我们将反复用定理~\ref{thm:6.12} 来描述极小曲面及其几乎共形参数化的
局部行为.

在二维情形, $d$ 次球面调和函数 $p_d(x)$ 可以写成某个常数乘以
$\mathrm{e}^{\ii\theta}(x_1+\ii x_2)^d$ 的实部, 其中 $\theta\in[0,2\pi)$.
复值向量函数 $f$ 的实部记为 $\operatorname{Re}f$.

\begin{lemma}
\label{lem:6.13}
设
\begin{equation}
v(x)=p_d(x)+q(x),
\label{eq:6.46}
\end{equation}
其中 $p_d$ 是 $d$ 次非平凡齐次调和多项式, 并且存在常数 $C$ 使
\begin{equation}
|q(x)|+|x||\nabla q(x)|\leq C|x|^{d+1};
\label{eq:6.47}
\end{equation}
则存在 $0$ 的邻域 $U$ 与 $C^1$ 微分同胚 $F:U\to D$, $F(0)=0$, 使得
$F(U\cap\{v=0\})$ 等于 $2d$ 条光滑弧
\begin{equation}
\alpha_j(t)=t\mathrm{e}^{\frac{j\pi\ii}{d}}
\label{eq:6.48}
\end{equation}
之并, 其中 $t\in[0,\epsilon]$, $j=1,\ldots,2d$.
\end{lemma}

\begin{proof}
由于处于二维情形, $p_d(x)$ 可写成某个常数乘以 $\mathrm{e}^{\ii\theta}z^d$
的实部, 其中 $z=x_1+\ii x_2$. 经一个旋转, 可设 $\theta=0$, 于是
\begin{equation}
v(z)=\operatorname{Re}\bigl(z^d+q(z)\bigr).
\label{eq:6.49}
\end{equation}
估计 \eqref{eq:6.47} 蕴涵映射
\begin{equation}
F(z)=z\bigl(1+z^{-d}q(z)\bigr)^{\frac1d}
\label{eq:6.50}
\end{equation}
在 $|z|\leq\frac{1}{2C}$ 上有定义. $F$ 的复导数为
\begin{align}
F_z(z)&=\bigl(1+z^{-d}q(z)\bigr)^{\frac1d}\notag\\
&\quad+\frac{z}{d}\bigl(1+z^{-d}q(z)\bigr)^{\frac{1-d}{d}}
\bigl(z^{-d}q_z(z)-dz^{-d-1}q(z)\bigr).
\label{eq:6.51}
\end{align}
再次应用 \eqref{eq:6.47} 可知 $F$ 是 $C^1$ 的且 $F_z(0)=1$. 由隐函数
定理即得所需的邻域与 $C^1$ 逆映射 $F^{-1}$.

与微分同胚 $F$ 复合, 便看出 $v(x)=\operatorname{Re}(F(z))^d$. 描述
\eqref{eq:6.48} 随之立即得到.
\end{proof}

\mtpara{Rado 定理}
定理~\ref{thm:6.12} (Bers) 的第一个应用是得到 Rado 的一个结果.

\begin{theorem}[Rado, \cite{Ra2}]
\label{thm:6.14}
设 $\Omega\subset\R^2$ 是凸子集, $\sigma\subset\R^3$ 是一条在 $\partial\Omega$
上图形化的简单闭曲线. 则任何以 $\partial\Sigma=\sigma$ 为边界的极小圆盘
$\Sigma\subset\R^3$ 必在 $\Omega$ 上图形化, 从而由极值原理唯一.
\end{theorem}

\begin{proof}
设 $\Sigma$ 是这样的一个非图形化极小圆盘. 则到 $(x_1,x_2)$-平面的投影
不可能是浸入. 因此存在某个 $x\in\Sigma$, 使得在 $x$ 处
$\nabla_\Sigma x_1$ 与 $\nabla_\Sigma x_2$ 线性相关. 换言之, 存在
$(a,b)\neq(0,0)$ 使
\begin{equation}
\nabla_\Sigma(ax_1+bx_2)(x)=0.
\label{eq:6.52}
\end{equation}
由命题~\ref{prop:1.7}, $ax_1+bx_2$ 在 $\Sigma$ 上调和 (因为它是坐标函数的
线性组合). 于是可以应用定理~\ref{thm:6.12}; 由 \eqref{eq:6.52} 与
引理~\ref{lem:6.13} 可知, 零点线
\begin{equation}
ax_1+bx_2=(ax_1+bx_2)(x)
\label{eq:6.53}
\end{equation}
在 $x$ 处有一个奇点, 至少有四条不同的曲线在该点相遇. 若其中两条零点
曲线再次相交, 则得到一条闭零点曲线, 它必须围成一个圆盘 (因为 $\Sigma$
是圆盘). 由极值原理, $ax_1+bx_2$ 在该圆盘上必为常数, 从而由唯一延拓在
$\Sigma$ 上为常数. 这将蕴涵 $\sigma=\partial\Sigma$ 含于 \eqref{eq:6.53}
给出的平面之内. 这是不可能的, 故所有这些曲线都通往边界而不再相交.

换言之, 由 \eqref{eq:6.53} 给出的 $\R^3$ 中的平面与 $\sigma$ 至少交于
四点. 然而, 由于 $\Omega\subset\R^2$ 是凸的, $\partial\Omega$ 与
\eqref{eq:6.53} 给出的直线恰好交于两点. 最后, 由于 $\sigma$ 在
$\partial\Omega$ 上图形化, $\sigma$ 与 \eqref{eq:6.53} 给出的 $\R^3$
中的平面恰好交于两点, 这就得到所需的矛盾.
\end{proof}

\mtpara{调和映射的标准形}
下面我们把定理~\ref{thm:6.12} 应用于调和映射 $u:D\to\R^3$, 以得到它的
一种标准形.

\begin{lemma}
\label{lem:6.15}
设 $u:D\to\R^3$ 是非常值调和映射, $u(0)=0$、$\nabla u(0)=0$. 则下述
局部渐近展开成立:
\begin{equation}
u(x)=\operatorname{Re}\sum_{j=d}^{2d-1}a_j(x_1+\ii x_2)^j+q(x),
\label{eq:6.54}
\end{equation}
其中 $d\geq2$, 每个 $a_i\in\C^3$, $a_d\neq0$, 且
\begin{equation}
|q(x)|+|x||\nabla q(x)|\leq C|x|^{2d}.
\label{eq:6.55}
\end{equation}
\end{lemma}

\begin{proof}
由于 $u$ 的各分量实际上是 $D$ 上的调和函数, 定理~\ref{thm:6.1} 蕴涵 $u$
在该点以某个有限阶 $d\geq2$ 消没. 记 $u^1$ 为该展开中的前 $d$ 项 (于是
$u^1$ 含有 $d$ 次到 $2d-1$ 次的项, 而 $d$ 次项非平凡). 于是可写
\begin{equation*}
u(x)=u^1(x)+q(x),
\end{equation*}
其中 $q(x)$ 满足 \eqref{eq:6.55}, 结论成立.
\end{proof}

\mtpara{分支点}
现在设 $u:D\to\R^3$ 是几乎共形调和映射 (于是其像为极小曲面). 我们在
第 4 章中已经看到, $u$ 在 (孤立的) 分支点之外是浸入, 而在分支点处
$u_{x_1}$ 与 $u_{x_2}$ 同时为零. 我们希望给出 $u$ 在这样的分支点附近的
描述. 我们的叙述大致遵循 Gulliver 在 \cite{Gu} 中的论证.

记 $z=x_1+\ii x_2$ 为 $D$ 上的复坐标. 由于 $u$ 几乎共形,
\begin{equation}
\langle u_z,u_z\rangle=0.
\label{eq:6.56}
\end{equation}
注意这给出两个方程 (即实部与虚部同时为零). 把它应用于局部表示
\eqref{eq:6.54}, 得到 $\langle a_d,a_d\rangle=0$ 与
$\langle a_d,a_{d+1}\rangle=0$. 前一个条件蕴涵 $a_d$ 的实部与虚部正交
且等长. 经一个旋转, 可设 $a_d=(a,-\ii a,0)$, 其中 $a>0$. 经伸缩, 可取
$a=1$, 于是 $a_d=(1,-\ii,0)$. 由于 $\langle a_d,a_{d+1}\rangle=0$,
可写 $a_{d+1}=(c,-\ii c,c')$. 把它代回 \eqref{eq:6.55}, 得到
\begin{equation}
u_1(z)+\ii u_2(z)=z^d+cz^{d+1}+G(z),
\label{eq:6.57}
\end{equation}
其中 $d\geq2$ 且
\begin{equation}
|z|\bigl|u_3(z)\bigr|+|G(z)|+|z|\bigl|\nabla G(z)\bigr|
\leq C|z|^{d+2}.
\label{eq:6.58}
\end{equation}
方程 \eqref{eq:6.57} 与 \eqref{eq:6.58} 给出了极小曲面在分支点附近的
一种标准形. 通过选取 $D$ 的一个好的重参数化, 可以得到更为典范的表示,
这将带来方便.

\begin{lemma}
\label{lem:6.16}
设 \eqref{eq:6.57} 与 \eqref{eq:6.58} 成立. 则存在 $0\in D$ 的邻域
$U$ 与 $V$, 以及 $C^1$ 微分同胚 $F:V\to U$, 使得对 $z\in V$,
\begin{equation}
u_1+\ii u_2=(F(z))^d,
\label{eq:6.59}
\end{equation}
且 $u_3=\phi(F(z))$, 其中
\begin{equation}
|\phi(z)|+|z||\nabla\phi(z)|+|z|^2|\nabla^2\phi(z)|\leq C|z|^{d+1}.
\label{eq:6.60}
\end{equation}
若 $u$ 实解析, 则 $F$ 与 $\phi$ 在 $0$ 之外实解析.
\end{lemma}

注意 \eqref{eq:6.60} 中的 $C$ 与 \eqref{eq:6.58} 中的 $C$ 不同. 可以取
$F$ 与 $\phi$ 为实解析这一事实, 要到研究假分支点时才会用到.

\begin{proof}
对 $G(z)$ 的估计 \eqref{eq:6.58} 蕴涵: 当 $z$ 充分小时,
\begin{equation}
|z|^{-d}|G(z)|<1.
\label{eq:6.61}
\end{equation}
因此可以定义
\begin{equation}
F(z)=z\bigl(1+z^{-d}G(z)\bigr)^{\frac1d}.
\label{eq:6.62}
\end{equation}
显然, 按此定义 \eqref{eq:6.59} 成立.

估计 \eqref{eq:6.58} 蕴涵 $F$ 是 $C^1$ 的且在原点处的导数为恒等映射.
因此由隐函数定理即得所需的邻域 $U$ 与 $V$ 以及 $C^1$ 逆映射 $F^{-1}$.
估计 \eqref{eq:6.60} 直接由 \eqref{eq:6.58} 得到.

若 $u$ 实解析, 则 $G$ 也实解析. 由 \eqref{eq:6.62} 立即看出 $F$ 在 $0$
之外实解析. 由于 $\phi$ 是实解析函数与 $F$ 的复合, 它在 $F$ 解析之处
(即 $0$ 之外) 实解析.
\end{proof}

像极小曲面局部地是平面 $x_3=0$ 上的多值图. 下面我们分析高度函数
$\phi$. 考虑到参数化 \eqref{eq:6.59}, 可以把引理~\ref{lem:6.16} 中
$\phi$ 的极小曲面方程写成
\begin{equation}
\dive\Bigl(\frac{\nabla\phi}{W}\Bigr)=0,
\label{eq:6.63}
\end{equation}
其中
\begin{equation}
W=\Bigl(1+\frac{|\nabla\phi|^2}{d^2|z|^{2d-2}}\Bigr)^{\frac12}.
\label{eq:6.64}
\end{equation}
当 $d=1$ 时, \eqref{eq:6.63} 就是经典的极小曲面方程 (见第 1 章). 当
$d\geq2$ 时, \eqref{eq:6.60} 蕴涵 $W$ 有界, 于是 \eqref{eq:6.63} 仍是
一致椭圆方程.

\begin{lemma}
\label{lem:6.17}
设 $d\geq1$ 固定, $\phi_1$ 与 $\phi_2$ 是 $D$ 上 \eqref{eq:6.63} 的解
($W$ 如 \eqref{eq:6.64}), $\phi_i(0)=0$ 且
$|\phi_i(z)|\leq C_0|z|^{d+1}$. 若 $\Phi=\phi_1-\phi_2$ 不恒为零, 则
存在整数 $n\geq d+1$ (事实上 $n\geq d+2$) 与渐近展开
\begin{equation}
\Phi(z)=\operatorname{Re}cz^n+\rho(z),
\label{eq:6.65}
\end{equation}
其中 $c\neq0$ 且
\begin{equation}
|\rho(z)|+|z||\nabla\rho(z)|\leq C|z|^{n+\epsilon}
\label{eq:6.66}
\end{equation}
对某个 $\epsilon>0$ 与 $C<\infty$ 成立.
\end{lemma}

\begin{proof}
正如在强极值原理 (引理~1.27) 的证明中那样, 我们断定 $\Phi$
满足形如
\begin{equation}
0=(a_{i,j}\Phi_{x_j})_{x_i}+b_i\Phi_{x_i}+c\Phi
\label{eq:6.67}
\end{equation}
的方程, 其中 $a_{ij},b_i,c$ 光滑, $a_{ij}$ 对称、一致椭圆且
$a_{ij}(0)=\delta_{ij}$.

现在可以对 $\Phi$ 应用定理~\ref{thm:6.12} 了. 我们断定: 若 $\Phi$ 不恒
为零, 则由 $\Phi(0)=0$ 可知存在形如 \eqref{eq:6.65} 的渐近展开.
\end{proof}

引理~\ref{lem:6.17} 在 $d\geq2$ 的情形将用于研究极小曲面的分支点. 在
后面的各节中, 当我们分析浸入极小曲面的相交时, 将把引理~\ref{lem:6.17}
应用于极小曲面方程 (即 $d=1$ 的情形).

现在设 $\zeta$ 是一个 $d$ 次单位根 (即 $\zeta^d=1$). 则
$\phi_\zeta(z)\equiv\phi(\zeta z)$ 也满足 \eqref{eq:6.63}. 此外, 极小
曲面的自交点正是 $\phi_\zeta$ 与 $\phi$ 相等之点 (这里允许 $\zeta$ 取
一切可能的值). 暂时固定 $\zeta$, 定义函数 $\Phi$ 为
\begin{equation}
\Phi(z)=\phi_\zeta(z)-\phi(z).
\label{eq:6.68}
\end{equation}
由引理~\ref{lem:6.17}, 若 $\Phi$ 不恒为零, 则得到形如
\begin{equation}
\Phi(z)=\operatorname{Re}cz^n+\rho(z)
\label{eq:6.69}
\end{equation}
的渐近展开, 其中 $c\neq0$ 且
\begin{equation}
|\rho(z)|+|z||\nabla\rho(z)|\leq C|z|^{n+\epsilon}
\label{eq:6.70}
\end{equation}
对某个 $\epsilon>0$ 成立.

完全仿照引理~\ref{lem:6.16} 的论证, 可以作一个在原点处共形的 $C^1$
坐标变换 $F$, 使得在某个球 $B_{r_0}$ 上可把 $\Phi$ 表示为
\begin{equation}
\Phi(z)=\operatorname{Re}F(z)^n.
\label{eq:6.71}
\end{equation}
最后, 表示 \eqref{eq:6.71} 使我们能够很好地描述 $u$ 在分支点附近的行为.

\begin{theorem}
\label{thm:6.18}
设 $u:D\to\R^3$ 是几乎共形调和映射 (于是其像为极小曲面). 则存在包含
原点的某个邻域 $V$, 使得或者 $u(V)$ 是浸入曲面, 或者存在简单 $C^1$ 弧
$\gamma_i:[0,1]\to D$, 满足 $\gamma_i(0)=0$、$|\gamma_i'(0)|=1$、
$\gamma_1'(0)\neq\gamma_2'(0)$、对一切 $t$ 有
$u(\gamma_1(t))=u(\gamma_2(t))$, 并且 $u$ 的像在 $\gamma_i$ 的像处的
切空间横截相交.
\end{theorem}

\begin{proof}
可设 $u$ 在 $0$ 处有分支点 (否则第一种情形平凡地成立). 由
引理~\ref{lem:6.16}, 存在邻域 $V$, 使得经坐标变换后有表示
$u_1+\ii u_2=z^n$ 与 $u_3(z)=\phi(z)$, 其中 \eqref{eq:6.60} 成立. 令
$\zeta=\mathrm{e}^{2\pi\ii/n}$, $\Phi=\phi_\zeta-\phi$.

若 $\Phi$ 恒为零, 则对一切 $z$ 有 $\phi(z)=\phi(\zeta z)$, 从而
$\phi(z)=\bar{\phi}(z^n)$, 其中 $\bar{\phi}$ 光滑. 于是可以取
$\bar{u}(z)=u(z^{1/n})$ 来重新参数化像 $u(V)$. $\bar{u}$ 的像是浸入
曲面, 且显然与 $u$ 的像重合.

否则, 若 $\Phi$ 不恒为零, 则在 $F$ 给出的坐标下有 $\Phi$ 的表示
\eqref{eq:6.71}. 在这些坐标下,
\begin{equation}
\Phi\bigl(t\mathrm{e}^{\frac{\pi\ii}{2n}}\bigr)
=\operatorname{Re}(\ii t^n)=0
\label{eq:6.72}
\end{equation}
对 $0\leq t\leq r_0$ 成立. 由于按定义 $\Phi$ 在 $\phi$ 与 $\phi_\zeta$
相等处为零, \eqref{eq:6.72} 蕴涵
\begin{equation}
\phi\bigl(t\mathrm{e}^{\frac{\pi\ii}{2n}}\bigr)
=\phi_\zeta\bigl(t\mathrm{e}^{\frac{\pi\ii}{2n}}\bigr)
=\phi\bigl(t\mathrm{e}^{\frac{5\pi\ii}{2n}}\bigr)
\label{eq:6.73}
\end{equation}
对 $0\leq t\leq r_0$ 成立. 因此, 在 $F$ 给出的坐标下, 曲线 $\gamma_i$
由 $B_{r_0}$ 中从原点出发、夹角为 $\frac{\pi}{2n}$ 与 $\frac{5\pi}{2n}$
的射线给出.

最后, 由于 $\{|\nabla\operatorname{Re}z^n|=0\}=\{0\}$, 表示
\eqref{eq:6.71} 蕴涵 $\nabla\Phi$ 沿 $\gamma_i$ 不消失 (原点除外). 因此,
$u$ 的像在 $\gamma_i$ 的像处的切空间横截相交.
\end{proof}

\begin{definition}
\label{def:6.19}
若 $0$ 是 $u$ 的分支点, 但存在包含原点的某个邻域 $V$ 使 $u(V)$ 是浸入
曲面, 则称 $0$ 为\newterm{false-branch}{假分支点}{false branch point}. 若存在包含原点的某个邻域 $V$ 与简单
$C^1$ 弧 $\gamma_i:[0,1]\to D$, 满足 $\gamma_i(0)=0$、$|\gamma_i'(0)|=1$、
$\gamma_1'(0)\neq\gamma_2'(0)$、对一切 $t$ 有
$u(\gamma_1(t))=u(\gamma_2(t))$, 并且 $u$ 的像在 $\gamma_i$ 的像处的切
空间横截相交, 则称 $0$ 为\newterm{true-branch}{真分支点}{true branch point}.
\end{definition}

假分支点来源于参数化, 而真分支点在像曲面的几何中是可见的. 出于这一
原因, 对面积最小化圆盘排除真分支点要容易得多. 我们将在下一节回到这个
问题.

% ================================================================
\section{真分支点的不存在性}
\label{sec:6.3}

本节我们将证明面积最小映射 $u:D\to\R^3$ 没有真分支点; 换言之, $u$ 的像
必是浸入曲面. 特别地, 若 $u$ 是 Plateau 问题的解 (见
定义~\ref{def:4.15}), 则 $u(D)$ 是浸入曲面. 真分支点已由 R. Osserman 在
\cite{Os1} 中排除. 这留下了 $u$ 具有假分支点的可能性 (此时问题出在
参数化而非曲面本身). 这将在下一节处理.

\begin{theorem}[Osserman, \cite{Os1}]
\label{thm:6.20}
若 $u:D\to\R^3$ 是几乎共形调和映射且在原点处有真分支点, 则存在映射
$v:D\to\R^3$, 使得 $u$ 与 $v$ 有相同的像, $u$ 与 $v$ 在 $\partial D$ 上
同胚, 并且 $v$ 关于某个紧支光滑变分不是能量泛函的驻定点.
\end{theorem}

作为直接的推论, 我们看到若 $u$ 有真分支点, 它就不可能是面积最小化的.

\begin{corollary}[Osserman, \cite{Os1}]
\label{cor:6.21}
若 $u:D\to\R^3$ 是 Plateau 问题的解, 则它不可能有任何真分支点.
\end{corollary}

\begin{proof}
若 $u$ 有真分支点, 则定理~\ref{thm:6.20} 将给出一个具有相同能量而非
极小化的新映射. 这导致矛盾.
\end{proof}

在定理~\ref{thm:6.20} 的证明中, 我们将利用真分支点的存在性, 通过剪贴
论证构造一个能量相同、但在某条曲线上发生折叠的映射. 由于圆盘上的驻定
映射只有孤立奇点, 新映射不可能是驻定的.

\begin{proof}[定理~\ref{thm:6.20} 的证明]
如所许诺, 我们将构造一个与 $u$ 有相同像的几乎共形映射 $v$, 使得
\begin{equation}
E_v=\Area_v=\Area_u=E_u,
\label{eq:6.74}
\end{equation}
且 $v$ 关于能量不是驻定点. 在此之前, 我们先构造一个与 $u$ 有相同像的
映射 $\bar{v}$, 它连续, 并且在一个零测集之外是浸入. Morrey 版本的
单值化定理保证我们可以重新参数化 $\bar{v}$ 以得到几乎共形映射 $v$.

设 $0$ 是 $u:D\to\R^3$ 的真分支点, $\gamma_i$ ($i=1,2$) 与 $V$ 由
定理~\ref{thm:6.18} 给出. 可设 $\partial V$ 光滑, 且每条曲线 $\gamma_i$
与 $\partial V$ 横截相交. 事实上, 把 $\gamma_i$ 重新参数化后, 可使
\begin{equation}
\gamma_i\cap\partial V=\gamma_i(\epsilon)
\label{eq:6.75}
\end{equation}
对某个 $\epsilon>0$ 成立.

选取同胚 $\bar{F}:\overline{D_\epsilon}\to\overline{V}$, 它在原点之外是
$C^2$ 的, 且对 $0\leq t\leq\epsilon$ 有 $\bar{F}(\ii t)=\gamma_1(t)$、
$\bar{F}(-\ii t)=\gamma_2(t)$.

下面我们构造映射 $G:\overline{D_\epsilon}\to\overline{D_\epsilon}$,
它具有下述性质. 第一, $G$ 在 $\partial D_\epsilon$ 上是恒等映射, 且在
$\partial D_\epsilon$ 的某个邻域内连续. 第二, 对
$-\frac{\epsilon}{2}\leq t\leq\frac{\epsilon}{2}$,
\begin{equation}
\lim_{s\to0_+}G(t+\ii s)=2\ii|t|
\label{eq:6.76}
\end{equation}
且
\begin{equation}
\lim_{s\to0_-}G(t+\ii s)=-2\ii|t|.
\label{eq:6.77}
\end{equation}
第三, 尽管 $G$ 在实轴两侧不连续, $u(\bar{F}(G))$ 却处处连续. 这是可以
做到的, 因为 \eqref{eq:6.76} 与 \eqref{eq:6.77} 蕴涵
$u(\bar{F}(G(t)))=u(\gamma_i(|t|))$. 第四, $u(\bar{F}(G))$ 是分段 $C^2$
的, 并且几乎处处是浸入.

显然, 构造映射 $G$ 有很大的自由度. 我们描述一种可能的构造. 选取不连续
映射 $G:\overline{D_\epsilon}\to\overline{D_\epsilon}$, 使得虚轴的负部与
正部分别映为 $-\ii\epsilon$ 与 $\ii\epsilon$, $-\frac{\epsilon}{2}$ 与
$\frac{\epsilon}{2}$ 映为原点, 不连续线段
$[-\frac{\epsilon}{2},0]$ 与 $[0,\frac{\epsilon}{2}]$ 按 \eqref{eq:6.76}
与 \eqref{eq:6.77} 映过去, $G$ 在 $\partial D_\epsilon$ 上为恒等映射且在
$\partial D_\epsilon$ 的某个邻域内连续, 并且 $G$ 在
\begin{equation}
D_\epsilon\setminus\Bigl(\bigl[-\tfrac{\epsilon}{2},0\bigr]
\cup\bigl[0,\tfrac{\epsilon}{2}\bigr]\cup[-\ii\epsilon,\ii\epsilon]\Bigr)
\label{eq:6.78}
\end{equation}
的每个连通分支上都是微分同胚. 映射 $G$ 的构造使得 $u(\bar{F}(G))$ 连续
且分段 $C^2$.

定义新映射 $\bar{v}$ 为
\begin{equation}
\begin{aligned}
\bar{v}(z)&=u\bigl(\bar{F}\bigl(G\bigl(\bar{F}^{-1}(z)\bigr)\bigr)\bigr)
\quad\text{当 }z\in V,\\
\bar{v}(z)&=u(z)\quad\text{当 }z\notin V.
\end{aligned}
\label{eq:6.79}
\end{equation}
映射 $\bar{v}$ 连续、分段 $C^2$, 并且在一个零测集之外是浸入.

Morrey 版本的单值化定理 (见注~\ref{rem:4.3}) 蕴涵存在 $W^2$ 中的同胚
$T:D\to D_\epsilon$ (于是其 Hessian 属于 $L^2$), 使 $v=\bar{v}\circ T$
几乎共形. 按构造, $u$ 与 $v$ 的像完全相同, 从而面积相等.

然而, 像极小曲面沿 $\gamma_i$ 的横截性蕴涵映射 $v$ 现在沿两整条线段都有
分支点. 因此 $v$ 关于能量不可能是驻定点, 因为由推论~\ref{cor:4.18},
驻定映射的分支点是孤立的.
\end{proof}

% ================================================================
\section{假分支点的不存在性}
\label{sec:6.4}

本节的主要结果定理~\ref{thm:6.23} 表明: 当 $\Omega$ 平均凸且
$u(\partial D)$ 嵌入时, Plateau 问题的解 $u:D\to\Omega\subset\R^3$ 没有
假分支点. 结合上一节 Osserman 的结果, 便知能量最小化映射 $u$ 是浸入.

这一结果最一般的形式属于 R. Gulliver, 即 \cite{Gu} 的定理 8.2 (另见
H. W. Alt \cite{Alt}). 回忆: 若 $u$ 在所有限制在 $\partial D$ 上为到
$\Gamma$ 的单调映射之映射中使能量最小, 则称 $u$ 是以 $\Gamma$ 为边界的
Plateau 问题的解 (见定义~\ref{def:4.15}).

\begin{theorem}[假分支点的不存在性, \cite{Gu}; 参看 \cite{Alt}]
\label{thm:6.22}
设 $M^3$ 是三维 $C^3$ Riemann 流形, $\Gamma\subset M$ 是分段 $C^1$
Jordan 曲线. 若 $u:D\to\R^3$ 是以 $\Gamma$ 为边界的 Plateau 问题的解,
则 $u$ 在 $D$ 上是浸入. 也就是说, $u$ 没有内部分支点.
\end{theorem}

由于我们的主要兴趣在于 $u(D)$ 的边界落在某个平均凸区域内的特殊情形,
我们将不再进一步讨论定理~\ref{thm:6.22}. 边界分支点的问题除了实解析
边界的情形 (由 Gulliver 与 F. Leslie 处理, 见 \cite{GuLe}) 之外尚未完全
解决. 另一方面, 对极小化流, 决定性的边界正则性结果由 Hardt 与 Simon 在
\cite{HaSi} 中证明 (关于极小曲面边界正则性研究的最新进展见 \cite{Wh2}).

本节的主要定理如下:

\begin{theorem}[假分支点的不存在性]
\label{thm:6.23}
设 $\Omega$ 是有界平均凸区域, 边界光滑, 并设
$u:\bar{D}\to\bar{\Omega}\subset\R^3$ 是如定义~\ref{def:4.15} 所述的
Plateau 问题的解. 若 $u(\partial D)\subset\partial\Omega$ 是嵌入的, 且
$u(D)\cap\Omega\neq\emptyset$, 则 $u(D)\subset\Omega$ 且 $u$ 没有假
分支点.
\end{theorem}

设 $u:D\to\Omega\subset\R^3$ 是 Plateau 问题的解, 于是 $u$ 几乎共形、
调和, 且 $u:\partial D\to\partial\Omega$ 单调. 在推论~\ref{cor:6.21}
(上一节) 中我们已经看到, 像 $u(D)$ 是浸入曲面.

由引理~\ref{lem:6.16}, 经一个在 $0$ 之外实解析的 $C^1$ 坐标变换后, 存在
$0\in D$ 的邻域 $V$ 与整数 $d<\infty$, 使得对 $z\in V$,
\begin{equation}
u_1+\ii u_2=z^d
\label{eq:6.80}
\end{equation}
且 $u_3=\phi(z)$, 其中
\begin{equation}
|\phi(z)|+|z||\nabla\phi(z)|+|z|^2|\nabla^2\phi(z)|\leq C|z|^{d+1}.
\label{eq:6.81}
\end{equation}
若 $d\geq2$, 则 $0$ 是 $u$ 的分支点. 此时, 若对一切 $z\in V$ 与任意
$d$ 次单位根 $\zeta_d$ 有
\begin{equation}
\phi(z)\equiv\phi(\zeta_dz),
\label{eq:6.82}
\end{equation}
则 $0$ 是假分支点. 由上一节的结果, 所有内部分支点都是假的.

当然, 若 $0$ 不是分支点, 则我们有 $d=1$ 时的 \eqref{eq:6.80} 与
\eqref{eq:6.81}.

在证明该定理之前, 我们需要一些预备引理.

\begin{lemma}
\label{lem:6.24}
对任意 $z\in\bar{D}$ 与紧集 $K\subset D$, 集合 $K\cap u^{-1}(u(z))$ 有限.
\end{lemma}

\begin{proof}
对任意点 $z_j\in K\cap u^{-1}(u(z))\subset D$, 表示 \eqref{eq:6.80}
蕴涵存在 $z_j$ 的邻域 $V_{z_j}$ 使
\begin{equation}
V_{z_j}\cap u^{-1}(u(z))=\{z_j\}.
\label{eq:6.83}
\end{equation}
由于 $u$ 在 $\bar{D}$ 上连续, 集合 $u^{-1}(u(z))$ 紧, 从而
$K\cap u^{-1}(u(z))$ 也紧. 因此存在有限个互不相同的点
$z_1,\ldots,z_k$, 使
\begin{equation}
K\cap u^{-1}(u(z))\subset\bigcup_{j=1}^{k}V_{z_j}.
\label{eq:6.84}
\end{equation}
结合 \eqref{eq:6.83} 与 \eqref{eq:6.84} 得
$K\cap u^{-1}(u(z))=\{z_j\}_{j=1,\ldots,k}$, 引理得证.
\end{proof}

下面我们转向极小曲面的边界含于某个平均凸区域边界内的情形. 这将是下一节
讨论的重点.

若该区域实际上是凸的, 则情况要简单一些, 如下例所示:

\begin{example}
\label{ex:6.25}
设 $u:\bar{D}\to\bar{B}_1\subset\R^3$ 是连续 (在 $\bar{D}$ 上) 调和
映射 (不必几乎共形), 满足 $u(\partial D)\subset\partial B_1$ 且
$u(D)\cap B_1\neq\emptyset$. 定义 $v:\bar{D}\to\R$ 为
$v(z)=1-|u(z)|^2$. 由于 $v\geq0$、$v(\partial D)=0$ 且 $\Delta v\leq0$,
强极值原理蕴涵 $v$ 在 $D$ 中为正. 换言之, $u(D)\subset B_1$. 此外,
Hopf 边界点引理蕴涵 $v$ 的法向导数在 $\partial D$ 上不为零. 若 $u$ 还是
几乎共形的且在 $\bar{D}$ 上 $C^1$, 则这一非零性蕴涵 $\partial D$ 的某个
邻域内没有分支点.
\end{example}

在下一个引理 (引理~\ref{lem:6.26}) 中, 我们将把此例推广到
$u:\bar{D}\to\bar{\Omega}\subset\R^3$ 是几乎共形调和映射、
$u(\partial D)\subset\partial\Omega$ 且 $\partial\Omega$ 平均凸的情形.
为此我们需要计算到平均凸区域边界的距离函数的 Hessian.

设 $\Omega\subset\R^3$ 是 $\R^3$ 中的平均凸区域, $\rho$ 表示到边界的
距离, 即
\begin{equation}
\rho(x)=\dist(\partial\Omega,x).
\label{eq:6.85}
\end{equation}
对 $y\in\partial\Omega$, 记 $\kappa_1(y)\leq\kappa_2(y)$ 为
$\partial\Omega$ 的主曲率. 由于边界 $\partial\Omega$ 平均凸, 有
\begin{equation}
0\leq\min_{y\in\partial\Omega}\bigl[\kappa_1(y)+\kappa_2(y)\bigr].
\label{eq:6.86}
\end{equation}
由于 $\partial\Omega$ 光滑且紧,
\begin{equation}
\max_{y\in\partial\Omega}\bigl[|\kappa_1(y)|+|\kappa_2(y)|\bigr]
\leq\bar{\kappa}<\infty.
\label{eq:6.87}
\end{equation}
结合 \eqref{eq:6.86} 与 \eqref{eq:6.87}, 对任意 $y\in\partial\Omega$ 与
$t<1/\bar{\kappa}$, 得到
\begin{equation}
\begin{aligned}
0&\leq\frac{\kappa_1(y)+\kappa_2(y)-2t\kappa_1(y)\kappa_2(y)}
{(1-\kappa_1(y)t)(1-\kappa_2(y)t)}\\
&=\frac{\kappa_1(y)}{1-\kappa_1(y)t}+\frac{\kappa_2(y)}{1-\kappa_2(y)t}.
\end{aligned}
\label{eq:6.88}
\end{equation}
若 $\kappa_1(y)\leq0$, 则 \eqref{eq:6.88} 是显然的. 若
$0<\kappa_1(y)$, 我们利用 $t<1/\bar{\kappa}$, 从而
\begin{equation}
2t\kappa_1(y)\kappa_2(y)\leq t\bigl(\kappa_1^2(y)+\kappa_2^2(y)\bigr)
<\kappa_1(y)+\kappa_2(y).
\label{eq:6.89}
\end{equation}
由 \eqref{eq:6.87}, 存在管状邻域
$\Omega_\delta=\{\rho<\delta\}\cap\Omega$, 使每个 $x\in\Omega_\delta$ 有
唯一的最近点 $x'\in\partial\Omega$. 换言之, 法指数映射是到
$\Omega_\delta$ 上的微分同胚. 因此我们可以用 Riccati 方程来计算距离函数
在 $\Omega_\delta$ 上到 $\partial\Omega$ 的 Hessian (\cite{GiTr} 第 355
页有此计算). 给定 $x\in\Omega$, 若
\begin{equation*}
\dist(\partial\Omega,x)=t<\delta,
\end{equation*}
则 $|\nabla\rho|=1$ ($\nabla\rho$ 是水平集 $\{\rho=t\}$ 的单位法向). 若
$x'\in\partial\Omega$ 是 $x$ 的最近点, 则 $\rho$ 的 Hessian 在 $x$ 处的
特征值为
\begin{equation}
\frac{-\kappa_1(x')}{1-\kappa_1(x')t},\quad
\frac{-\kappa_2(x')}{1-\kappa_2(x')t},\quad 0.
\label{eq:6.90}
\end{equation}
给定任意 $C^2$ 函数 $f:\R\to\R$, $f\circ\rho$ 的 Hessian 为
\begin{equation}
\bigl(f(\rho)\bigr)_{x_ix_j}
=f''(\rho)\rho_{x_i}\rho_{x_j}+f'(\rho)\rho_{x_ix_j}.
\label{eq:6.91}
\end{equation}
若对某常数 $a>0$ 取 $f(t)=2at-t^2$, 于是 $f\circ\rho=2a\rho-\rho^2$, 则
$f'(t)=2a-2t$、$f''(t)=-2$. 结合 \eqref{eq:6.90} 与 \eqref{eq:6.91}, 若
$x'\in\partial\Omega$ 是 $x$ 的最近点, 则 $2a\rho-\rho^2$ 的 Hessian 在
$x$ 处的特征值为
\begin{equation}
-2(a-t)\frac{\kappa_1(x')}{1-\kappa_1(x')t},\quad
-2(a-t)\frac{\kappa_2(x')}{1-\kappa_2(x')t},\quad -2.
\label{eq:6.92}
\end{equation}

\begin{lemma}
\label{lem:6.26}
设 $\Omega$ 是有界平均凸区域, 边界光滑, 并设
$u:\bar{D}\to\bar{\Omega}\subset\R^3$ 连续、几乎共形且调和. 若
$u(\partial D)\subset\partial\Omega$ 且 $u(D)\cap\Omega\neq\emptyset$,
则 $u(D)\subset\Omega$, 且 $|\nabla u|$ 在 $\partial D$ 上不为零.
\end{lemma}

注意这蕴涵 $u(\bar{D})$ 与边界横截相交.

\begin{proof}
选取充分小的 $a>0$ 使
\begin{equation}
a\bar{\kappa}<\frac12;
\label{eq:6.93}
\end{equation}
于是当 $t<\delta_1=\min\{\delta,\frac{1}{2\bar{\kappa}}\}$ 时,
\begin{equation}
(a-t)\frac{\bar{\kappa}}{1-\bar{\kappa}t}<1.
\label{eq:6.94}
\end{equation}
置
\begin{equation}
f(t)=2at-t^2
\label{eq:6.95}
\end{equation}
以及
\begin{equation}
v(z)=f\circ\rho(u(z)).
\label{eq:6.96}
\end{equation}
结合 \eqref{eq:6.92} 与 \eqref{eq:6.94}, 得到: 对 $\rho(y)\leq\delta_1$
的 $\bar{y}$, $\mathrm{Hess}_{f\circ\rho}(\bar{y})$ 的特征值满足
\begin{equation}
-2<-2(a-t)\frac{\kappa_2(y)}{1-\kappa_2(y)t}
\leq-2(a-t)\frac{\kappa_1(y)}{1-\kappa_1(y)t},
\label{eq:6.97}
\end{equation}
其中 $y\in\partial\Omega$ 是 $\bar{y}$ 的最近点.

我们将证明: 若 $u(z)$ 位于 $\partial\Omega$ 的某个管状邻域内, 则
$v(z)$ 超调和. 由此同样可推出 $u(D)\subset\Omega$, 且 $u$ 的法向导数在
边界上不为零.

设 $z=x_1+\ii x_2$ 为 $D$ 的坐标, $y_1,y_2,y_3$ 为 $\R^3$ 的坐标.
利用链式法则与 $\Delta u=0$, 得到
\begin{align}
\Delta v(z)&=v_{x_ix_i}(z)
=(f\circ\rho)_{y_jy_k}(u(z))(u_j)_{x_i}(z)(u_k)_{x_i}(z)\notag\\
&\quad+(f\circ\rho)_{y_j}(u(z))(u_j)_{x_ix_i}(z)\notag\\
&=(f\circ\rho)_{y_jy_k}(u(z))(u_j)_{x_i}(z)(u_k)_{x_i}(z).
\label{eq:6.98}
\end{align}
若以 $\Pi_z$ 记 $u_{x_1}(z)$ 与 $u_{x_2}(z)$ 张成的二平面, 则
\eqref{eq:6.98} 结合 $u$ 的几乎共形性给出
\begin{align}
\Delta v(z)&=\mathrm{Hess}_{f\circ\rho}(u(z))\bigl(u_{x_i}(z),u_{x_i}(z)\bigr)
\notag\\
&=|u_{x_1}(z)|^2\operatorname{Tr}
\mathrm{Hess}_{f\circ\rho}(u(z))\Bigr|_{\Pi_z}.
\label{eq:6.99}
\end{align}
现在设 $\rho(u(z))=t\leq\delta_1$, 于是 \eqref{eq:6.97} 适用, 从而
\begin{equation}
\begin{aligned}
&\max_{\text{二平面 }\Pi}\operatorname{Tr}
\mathrm{Hess}_{f\circ\rho}(u(z))\Bigr|_{\Pi}\\
&\quad=-2(a-t)\Bigl(\frac{\kappa_1(y)}{1-\kappa_1(y)t}
+\frac{\kappa_2(y)}{1-\kappa_2(y)t}\Bigr)\leq0,
\end{aligned}
\label{eq:6.100}
\end{equation}
其中最后一个不等式由 \eqref{eq:6.88} 得到, 而 $y$ 是 $\partial\Omega$
中距 $u(z)$ 最近的点.

结合 \eqref{eq:6.99} 与 \eqref{eq:6.100} 可知, 只要
\begin{equation*}
\rho(u(z))\leq\min\{\delta_1,a\},
\end{equation*}
$v$ 就是超调和的. 由于 $v$ 非负、在边界上为零, 且
$u(D)\cap\Omega\neq\emptyset$, 由极值原理它必在 $D$ 上严格为正. 按 $v$
的定义, 这意味着 $u(D)\subset\Omega$. 最后, 由于 $v$ 在边界每一点取到
最小值, Hopf 边界点引理蕴涵法向导数 $\frac{dv}{dn}$ 在边界上严格为负.
把它与链式法则结合, 得到在弱意义下
\begin{equation}
0<\frac{dv}{dn}=(f\circ\rho)_{y_j}(u)\frac{du_j}{dn}
=2a\Bigl\langle\nabla\rho(u),\frac{du}{dn}\Bigr\rangle,
\label{eq:6.101}
\end{equation}
从而 $\nabla u$ 在边界上不为零.
\end{proof}

作为一个推论, 我们看到 $u$ 在边界上的限制是嵌入, 且 $u$ 至多有有限个
内部分支点.

\begin{corollary}
\label{cor:6.27}
设 $\Omega$ 是有界平均凸区域, $\partial\Omega$ 光滑, 并设
$u:\bar{D}\to\bar{\Omega}\subset\R^3$ 在 $\bar{D}$ 上 $C^1$、几乎共形
且调和. 以 $\mathcal{B}_0$ 记 $u$ 的分支点, 并定义
$\mathcal{B}=u^{-1}(u(\mathcal{B}_0))$.

若 $u(\partial D)\subset\partial\Omega$ 且 $u(D)\cap\Omega\neq\emptyset$,
则 $\mathcal{B}$ 有限, 且存在紧集 $K\subset D$ 使
$\mathcal{B}\subset K$. 特别地, $u$ 在 $\partial D$ 上的限制是嵌入.
\end{corollary}

\begin{proof}
上一引理 (见 \eqref{eq:6.101}) 给出 $|\nabla u|>0$ 于 $\partial D$ 上,
因此 (因为 $u$ 几乎共形) $D$ 的边界上没有分支点. 由于
$u\in C^1(\bar{D})$, 这蕴涵 $\nabla u$ 在边界的某个邻域内不为零 (于是
边界附近没有分支点). 因此分支点紧含于 $D$ 之中. 由于局部描述
\eqref{eq:6.80} 蕴涵它们在 $D$ 中是孤立的, 分支点只能有有限多个. 令
\begin{equation}
\mathcal{B}_0=\{z_1,\ldots,z_k\}
\label{eq:6.102}
\end{equation}
表示这些分支点, 并令
\begin{equation}
\mathcal{B}=u^{-1}\bigl(u(\{z_1,\ldots,z_k\})\bigr).
\label{eq:6.103}
\end{equation}
上一引理蕴涵对每个 $j$ 有 $u^{-1}(u(z_j))\subset D$. 由于 $u$ 连续,
这蕴涵 $u^{-1}(u(z_j))$ 紧含于 $D$. 因此, 引理~\ref{lem:6.24} 蕴涵
$u^{-1}(u(z_j))$ 是有限集. 对每个 $j=1,\ldots,k$ 作此讨论, 便知
$\mathcal{B}$ 是紧含于 $D$ 的有限点集.
\end{proof}

现在我们准备证明本节的主要定理.

\begin{proof}[定理~\ref{thm:6.23} 的证明]
我们将假设 $0$ 是 $u$ 的 $d\geq2$ 阶假分支点, 且 $u(D)$ 不包含在任何
平面内, 并由此导出矛盾.

先证明存在过 $u(0)$ 的某个二平面 $\Pi$, 使 $\Pi$ 与 $u(D)$ 横截相交,
且
\begin{equation}
\Pi\cap u(\mathcal{B})=\{u(0)\}.
\label{eq:6.104}
\end{equation}
以 $\tilde{\Pi}$ 记 $\R^3$ 中过 $u(0)$ 且满足 \eqref{eq:6.104} 的二
平面之集.

由于 $u(\mathcal{B})$ 是有限点集, 只有有限多条过 $u(0)$ 的直线与
$u(\mathcal{B})\setminus\{u(0)\}$ 相交. 每条这样的直线含于过 $u(0)$ 的
一个单参数二平面族中, 故 $\tilde{\Pi}$ 是满测度的开集.

定义映射
$P:\bS^1\times(-\frac{\pi}{4},\frac{\pi}{4})\times\R^2\to\R^3$ 为
\begin{equation}
\begin{aligned}
P\bigl(\theta_1,\theta_2,(a_1,a_2)\bigr)
&=u(0)+a_1(-\sin\theta_1,\cos\theta_1,0)\\
&\quad+a_2(\cos\theta_1\cos\theta_2,\sin\theta_1\cos\theta_2,
-\sin\theta_2),
\end{aligned}
\label{eq:6.105}
\end{equation}
于是对每个 $(\theta_1,\theta_2)$ 得到一个仿射映射, 其像为过 $u(0)$ 的
二平面. 由于 $P$ 的微分在每点处满射 (作为到 $\R^3$ 的映射), Sard 定理
的参数形式适用. 因此, 对任意紧浸入子流形 $N\subset\R^3$, 使得
\begin{equation}
P(\theta_1,\theta_2,\R^2)
\label{eq:6.106}
\end{equation}
与 $N$ 横截相交的 $(\theta_1,\theta_2)$ 之集是满测度的. 特别地, 可以找到
一个二平面 $\Pi\in\tilde{\Pi}$ 与 $u(D)$ 横截相交.

经一个固定 $u(0)$ 的 $\R^3$ 旋转与一个平移, 可设 $u(0)=0$ 且 $\Pi$ 由
$y_1=0$ 给出. $\Pi$ 与 $u(D)$ 横截相交这一事实蕴涵
\begin{equation}
\nabla_{u(D)}y_1\neq0,
\label{eq:6.107}
\end{equation}
从而零点集 $\{y_1=0\}\cap u(D)$ 是一族紧浸入弧. 以 $\tilde{\gamma}$ 记
$\{y_1=0\}\cap u(D)$ 中含 $u(0)$ 的连通分支. 可写
\begin{equation}
\tilde{\gamma}=\bigcup_{j=1}^{\ell}\gamma_j,
\label{eq:6.108}
\end{equation}
其中每条 $\gamma_j$ 是紧浸入曲线, 或为闭曲线, 或有两个端点. 设
$u(0)\in\gamma_1$, 从而
$0\in u^{-1}(\gamma_1)\subset\{u_1=0\}$.

此外, \eqref{eq:6.104} 与 \eqref{eq:6.107} 蕴涵 $u^{-1}(\gamma_1)$ 是
一条在 $0$ 之外光滑的连通曲线. 由于 $0$ 是 $u$ 的 $d\geq2$ 阶分支点,
我们有 $u$ 的描述 \eqref{eq:6.80}. 因此, 在 $0$ 的某个邻域内,
$u^{-1}(\gamma_1)$ 是 $2d$ 条互不相交的弧之并.

下面我们断言这 $2d$ 条弧再也不会相交. 若其中两条再次相交, 则调和函数
$u_1$ 将有一个有界的零点域. 此时由极值原理, $u_1$ 必恒为零, 从而
$u(D)$ 将含于平面 $\{y_1=0\}$ 之内. 因此这 $2d$ 条弧一直互不相交直到
边界. 每条弧起于 $0$ 而终于 $\partial D$ (这里我们再次用了
引理~\ref{lem:6.26}). 这些曲线在 $\partial D$ 中的 $2d$ 个端点被 $u$
映为 $\gamma_1$ 的两个端点. 我们断定有两条不同的弧 $\alpha_1$ 与
$\alpha_2$ 满足
\begin{equation}
u(\alpha_i(0))=u(0)\quad\text{且}\quad
u(\alpha_1(1))=u(\alpha_2(1)),
\label{eq:6.109}
\end{equation}
其中 $u(\alpha_1(1))\in\partial u(D)$. 以 $\sigma$ 记 $\partial D$ 中
介于 $\alpha_1(1)$ 与 $\alpha_2(1)$ 之间的部分. 由于 $u$ 在边界上单调,
且曲线 $u(\partial D)$ 是嵌入的, $u^{-1}(u(\alpha_1(1)))$ 连通, 从而
$u(\sigma)=u(\alpha_1(1))$. 因此, 极值原理蕴涵 $u_1$ 在由 $\alpha_1$、
$\alpha_2$ 与 $\sigma$ 围成的区域上为零. 由唯一延拓, $u_1$ 在 $D$ 上恒
为零, 这就得到所需的矛盾.
\end{proof}

% ================================================================
\section{Plateau 问题的嵌入解}
\label{sec:6.5}

本节我们将看到, 在某些情形下 Plateau 问题的解必是嵌入的. 遵循 Meeks
与 Yau 的做法, 我们将看到: 若边界曲线是嵌入的, 且位于某个光滑凸集的
边界上 (并且它在该凸集中零伦), 则极小化的解是嵌入的. 在上一节中我们已经看到, 此时解必是浸入的. 注意: 要证明这样的
嵌入性定理, 对边界曲线作某种限制当然是必要的. 例如, 若边界曲线打结
(比如三叶结), 则它不可能被任何嵌入圆盘张成 (无论是否极小).

\begin{theorem}[Meeks--Yau, \cite{MeY1}]
\label{thm:6.28}
设 $M^3$ 是紧 Riemann 三维流形, 其边界平均凸, $\gamma$ 是 $\partial M$
中在 $M$ 内零伦的简单闭曲线; 则 $\gamma$ 被一个面积最小圆盘所界定, 且
任何这样的面积最小圆盘都是真嵌入的.
\end{theorem}

定理~\ref{thm:6.28} 含于他们的结果之中 (更多结果见 \cite{MeY1} 与
\cite{MeY2} 及其中的参考文献).

我们将讨论上述定理在 $M$ 为 $\R^3$ 中紧平均凸区域时的特殊情形.

\begin{theorem}[Meeks--Yau, \cite{MeY1}]
\label{thm:6.29}
设 $\Omega$ 是 $\R^3$ 中有界平均凸区域, 边界光滑, $\Gamma\subset\partial\Omega$
是 $C^1$ 简单闭曲线. 若 $\Gamma$ 在 $\Omega$ 中零伦, 则 Plateau 问题的解
$u:D\to\Omega$ 是真嵌入.
\end{theorem}

注意 Rado 定理, 即定理~\ref{thm:6.14}, 可看作定理~\ref{thm:6.29} 的
特殊情形. 此时边界曲线位于某个凸柱面的边界上. 一般地, 若一条曲线位于
其凸包的边界上, 则称之为\newterm{extremal}{极端的}{extremal}. 在 Meeks 与 Yau 的工作之前,
Gulliver 与 J. Spruck 的工作 \cite{GuSp} 已对 $\R^3$ 中总曲率较小的
极端边界曲线证明了嵌入性. 其后, Almgren 与 Simon \cite{AmSi}, 以及
Tomi 与 A. J. Tromba \cite{ToTr}, 对 $\R^3$ 中的极端边界曲线证明了
某个嵌入解的存在性 (但未必是\secref{sec:4.3}中给出的 Douglas--Rado 解).

由第 4 章与本章前几节的结果, 存在 Plateau 问题的解 $u$, 它是光滑真浸入,
且其在 $\partial D$ 上的限制是嵌入. 在本节余下的部分, 总假定 $u$ 具有
这些性质.

\mtpara{边界附近的嵌入性}
下述引理表明 $u$ 在边界附近是嵌入:

\begin{lemma}
\label{lem:6.30}
对如上所述的 $u$, 存在 $\delta,\bar{\delta}>0$, 使得
$u:D\setminus D_{1-\delta}\to\Omega$ 是嵌入, 且 $u(D)$ 在
$\partial\Omega$ 的 $\bar{\delta}$ 邻域内是嵌入的.
\end{lemma}

\begin{proof}
我们将证明: 只要 $\delta$ 充分小, 对 $D\setminus D_{1-\delta}$ 中的
任意 $w_1,w_2$ 都有 $u(w_1)\neq u(w_2)$.

由于 $u$ 是浸入且在 $\bar{D}$ 上 $C^1$, 存在 $\epsilon_0>0$, 使得对任意
$z_0\in D$,
\begin{equation}
u\bigl(D_{\epsilon_0}(z_0)\cap D\bigr)
\label{eq:6.110}
\end{equation}
是嵌入曲面. 因此可以假设 $|w_1-w_2|>\epsilon_0$.

由于 $u$ 在 $\bar{D}$ 上 $C^1$, 存在常数 $C_0$, 使得对任意
$z_1,z_2\in\bar{D}$,
\begin{equation}
\bigl|u(z_1)-u(z_2)\bigr|\leq C_0|z_1-z_2|.
\label{eq:6.111}
\end{equation}
$u$ 在 $\partial D$ 上的限制是 $C^1$ 嵌入这一事实蕴涵存在常数 $C_1>0$,
使得对任意 $z_1,z_2\in\partial D$,
\begin{equation}
C_1|z_1-z_2|\leq\bigl|u(z_1)-u(z_2)\bigr|.
\label{eq:6.112}
\end{equation}
取 $z_1,z_2\in\partial D$ 使 $|w_i-z_i|\leq\delta$. 由三角不等式以及
\eqref{eq:6.111} 与 \eqref{eq:6.112},
\begin{equation}
\begin{aligned}
|u(w_1)-u(w_2)|&\geq|u(z_1)-u(z_2)|-2C_0\delta
\geq C_1|z_1-z_2|-2C_0\delta\\
&\geq C_1(\epsilon_0-2\delta)-2C_0\delta.
\end{aligned}
\label{eq:6.113}
\end{equation}
式 \eqref{eq:6.113} 蕴涵: 只要
$\delta<C_1\epsilon_0/(2(C_0+C_1))$, $u$ 在 $D\setminus D_{1-\delta}$ 上
的限制就是嵌入.

由上一节的引理~\ref{lem:6.26}, 若
\begin{equation*}
v=f\circ\rho\circ u,
\end{equation*}
其中对某 $a>0$ 取 $f(t)=2at-t^2$, 而 $\rho$ 是到 $\partial\Omega$ 的
距离, 则当 $v(z)\leq\delta_0$ 时 $v$ 超调和, 其中 $\delta_0>0$; 并且
存在 $\beta>0$, 使得对一切 $z\in\partial D$,
\begin{equation}
\frac{\partial v}{\partial n}(z)>\beta.
\label{eq:6.114}
\end{equation}
由于 $u$ 在 $\bar{D}$ 上 $C^1$, 存在 $\delta_3>0$, 使当
$|z|>1-\delta_3$ 时 \eqref{eq:6.114} 成立. 因此存在
$\delta_1,\delta_2>0$ 与 $U\subset D$, 使 $D_{1-\delta_1}\subset U$ 且
$v(\partial U)=\delta_2<\delta_0$. 如上一节那样, 极值原理蕴涵
$v(U)\geq\delta_2$. 因此存在某个 $\bar{\delta}>0$, 使
\begin{equation}
u^{-1}\bigl(\{x\in\Omega\mid\dist(x,\partial\Omega)<\bar{\delta}\}\bigr)
\subset D\setminus D_{1-\delta}.
\label{eq:6.115}
\end{equation}
第二个结论现在由第一个结论推出.
\end{proof}

\mtpara{密度与双重点集}
回忆: 对点 $x\in\Omega$, $u(D)$ 在 $x$ 处的密度由
\begin{equation}
\Theta_x=\lim_{s\to0}
\frac{\Vol\bigl(B_s(x)\cap u(D)\bigr)}{\pi s^2}
\label{eq:6.116}
\end{equation}
给出. 若 $x\in u(D)$, 则在任何充分小的球 $B_\epsilon(x)$ 内 (由于
$u(D)$ 是浸入子流形), $B_\epsilon(x)\cap u(D)$ 由若干个在 $x$ 处相交的
极小曲面组成. 由单调性, 如推论~\ref{cor:1.13} 与 \ref{cor:1.14} 那样,
我们看到 $\Theta_x$ 等于 $x$ 的原像个数. 也就是说,
\begin{equation}
\Theta_x=\bigl|u^{-1}(x)\bigr|.
\label{eq:6.117}
\end{equation}
显然, 若 $x\notin u(D)$, 则 $\Theta_x=0$. 两种情形下 \eqref{eq:6.117}
都成立.

于是, \newterm{double-point}{双重点集}{double point set} $\mathcal{D}$ 就是使
\begin{equation}
\Theta_x\geq2
\label{eq:6.118}
\end{equation}
的 $x\in\Omega$ 之集. 由 \eqref{eq:6.117}, $\mathcal{D}$ 是使
$u(D)\setminus\mathcal{D}$ 嵌入的最小集合. 引理~\ref{lem:6.30} 蕴涵
\begin{equation}
\mathcal{D}\subset
\{x\in\Omega\mid\dist(x,\partial\Omega)\geq\bar{\delta}\}.
\label{eq:6.119}
\end{equation}
利用单调性、命题~\ref{prop:1.12} 与引理~\ref{lem:6.30},
\eqref{eq:6.117} 给出引理~\ref{lem:6.24} 的一个定量形式. 也就是说,
存在
\begin{equation}
C_2=C_2\bigl(\Vol(u(D)),\bar{\delta}\bigr)<\infty
\label{eq:6.120}
\end{equation}
使 $u(D)$ 中每一点的重数至多为 $C_2$.

\begin{lemma}
\label{lem:6.31}
对如上所述的 $u$, 双重点集 $\mathcal{D}\subset u(D)$ 是紧的.
\end{lemma}

\begin{proof}
引理~\ref{lem:6.30} 表明 $\mathcal{D}$ 含于 $u(D)$ 的内部, 于是可以对
$u(D)$ 在 $\mathcal{D}$ 的每一点附近应用单调性.

由推论~\ref{cor:1.14}, 密度是上半连续的, 且 $\mathcal{D}$ 是闭的. 由于
它还有界, 故 $\mathcal{D}$ 紧.
\end{proof}

\mtpara{非横截双重点}
利用椭圆方程解的局部描述, 我们将证明非横截自交之集
$\tilde{\mathcal{D}}\subset\mathcal{D}$ 是有限点集. 按定义,
$x\in\tilde{\mathcal{D}}$ 当且仅当 $u^{-1}(x)$ 中存在点 $z_1\neq z_2$,
以及 $z_1$ 的邻域 $V_1$ 与 $z_2$ 的邻域 $V_2$, 使得 $u(V_1)$ 与
$u(V_2)$ 在 $x$ 处相切相交. 此外, 可设每个 $u(V_i)$ 是某个函数
$\phi_i$ 在其于 $x$ 处的切平面中某个球 $B_\epsilon$ 上的图.

现在可以对 $\Phi=\phi_1-\phi_2$ 应用引理~\ref{lem:6.17} (由唯一延拓,
它不恒为零). 也就是说, 存在整数 $n\geq2$ 与渐近展开
\begin{equation}
\Phi(z)=\operatorname{Re}cz^n+\rho(z),
\label{eq:6.121}
\end{equation}
其中 $c\neq0$ 且
\begin{equation}
|\rho(z)|+|z||\nabla\rho(z)|\leq C|z|^{n+\alpha}
\label{eq:6.122}
\end{equation}
对某个 $\alpha>0$ 成立. 由引理~\ref{lem:6.13} 的一个小变形立即推出
$u(V_1)\cap u(V_2)$ 同胚于在 $x$ 处相交的 $2n$ 条弧.

\begin{lemma}
\label{lem:6.32}
对如上所述的 $u$ 与 $\mathcal{D}$, 以
$\tilde{\mathcal{D}}\subset\mathcal{D}$ 记具有非横截自交的双重点之集;
则 $\tilde{\mathcal{D}}$ 是有限点集.
\end{lemma}

\begin{proof}
设 $x\in\tilde{\mathcal{D}}$, 且
\begin{equation}
u^{-1}(x)=\{z_1,\ldots,z_k\}\subset D.
\label{eq:6.123}
\end{equation}
我们只讨论所有非横截叶瓣具有同一平面的情形. 一般情形的修改是清楚的.

由于 $u$ 是 $C^1$ 浸入, 存在 $z_i$ 的邻域 $V_i$, 使得 $u(V_i)$ 在
$P_i$ 上图形化, 其中 $P_i$ 是 $u(V_i)$ 在 $x$ 处的切平面, $i=1,\ldots,k$.

由于 $u$ 是固有的, 存在某个 $\epsilon_0>0$ 使
\begin{equation}
u^{-1}\bigl(B_{\epsilon_0}(x)\bigr)\subset\bigcup_{j=1}^{k}V_j.
\label{eq:6.124}
\end{equation}
按 $\tilde{\mathcal{D}}$ 的定义, 可设 $P_1,\ldots,P_\ell$ 在相差定向的
意义下都相同, 且 $\pm P_1$ 与 $P_{\ell+1},\ldots,P_k$ 之间有确定的
间隔. 因此存在某个 $\epsilon_1>0$, 使得对任意 $j_0\leq\ell<j_1$,
$B_{\epsilon_1}(x)\cap u(V_{j_0})$ 与 $B_{\epsilon_1}(x)\cap u(V_{j_1})$
横截相交.

最后, 本引理之前的讨论对每个 $j_0\neq j_1\leq\ell$ 给出某个
$\epsilon_{j_0,j_1}>0$, 使得 $B_{\epsilon_{j_0,j_1}}(x)\cap u(V_{j_0})$
与 $B_{\epsilon_{j_0,j_1}}(x)\cap u(V_{j_1})$ 在 $x$ 之外横截相交. 若取
$\epsilon>0$ 为 $\epsilon_0$、$\epsilon_1$ 以及所有
$\epsilon_{j_0,j_1}$ 中的最小值, 则
\begin{equation}
B_\epsilon(x)\cap\tilde{\mathcal{D}}=\{x\},
\label{eq:6.125}
\end{equation}
从而 $\tilde{\mathcal{D}}$ 是孤立的点集.

还需证明 $\tilde{\mathcal{D}}$ 紧. 为此设 $x_j\in\tilde{\mathcal{D}}$,
并设 $z_j^i\in u^{-1}(x_j)$ ($i=1,2$) 是使得它们的邻域在 $x_j$ 处相切
相交的点. 由于 $u$ 是 $C^1$ 浸入, 存在 $\epsilon_0>0$ 使
\begin{equation}
\bigl|z_j^1-z_j^2\bigr|\geq\epsilon_0.
\label{eq:6.126}
\end{equation}
由引理~\ref{lem:6.30}, $z_j^i$ 紧含于 $\Omega$ 之中, 从而有极限点
$z^1,z^2\in D$, 它们也满足 \eqref{eq:6.126} (于是 $z_1\neq z_2$).

由于 $u$ 是 $C^1$ 的, 由 $u(z_j^1)=u(z_j^2)$ 与
$|N(z_j^1)|=|N(z_j^2)|$ 对每个 $j$ 成立, 可知 $u(z^1)=u(z^2)$ 且
$|N(z^1)|=|N(z^2)|$; 又因 $z_1\neq z_2$, 故
$u(z^1)=u(z^2)\in\tilde{\mathcal{D}}$. 由 $z_j^1\to z_1$ 与 $u$ 的
连续性可知 $x_j\to u(z_1)$, 从而 $\tilde{\mathcal{D}}$ 紧.
\end{proof}

由上一引理可知, $\mathcal{D}$ 是一族紧浸入曲线, 它们在有限点集
$\tilde{\mathcal{D}}$ 处分叉.

\begin{corollary}
\label{cor:6.33}
对如上所述的 $u$ 与 $\mathcal{D}$, 可写
\begin{equation}
\mathcal{D}=\bigcup_{j=1}^{n}\eta_j,
\label{eq:6.127}
\end{equation}
其中每条 $\eta_j$ 是紧浸入曲线.
\end{corollary}

\begin{proof}
隐函数定理蕴涵 $\mathcal{D}\setminus\tilde{\mathcal{D}}$ 是 $C^1$ 浸入
曲线之并. 因此, 由引理~\ref{lem:6.32} 知 $\tilde{\mathcal{D}}$ 是有限
点集, 由引理~\ref{lem:6.31} 知 $\mathcal{D}$ 紧, 推论随之成立.
\end{proof}

\mtpara{折叠曲线}
最后一个预备知识是 Meeks 与 Yau 所用的折叠曲线的概念.

\begin{definition}
\label{def:6.34}
设 $f:D_r\to\R^3$ 是 Lipschitz 映射, 且 $f$ 在 $x_1\geq0$ 或
$x_1\leq0$ 上的限制是直到边界的 $C^1$ 浸入. 若对每个 $x_2$ 下述两种
可能之一出现, 则称 $f$ 沿 $x_2$-轴有一条\newterm{folding-curve}{折叠曲线}{folding curve}:

第一, $f_{x_2}(0,x_2)$ 与
$\lim_{x_1\to0_+}f_{x_1}(x_1,x_2)$ 张成的平面与 $f_{x_2}(0,x_2)$ 与
$\lim_{x_1\to0_-}f_{x_1}(x_1,x_2)$ 张成的平面横截. 第二, 向量
$\lim_{x_1\to0_-}f_{x_1}(x_1,x_2)$ 是
$\lim_{x_1\to0_+}f_{x_1}(x_1,x_2)$ 的负倍数.
\end{definition}

折叠曲线的要点在于: 如以下引理所示, 它们不可能出现在面积最小映射中:

\begin{lemma}
\label{lem:6.35}
若 $f$ 是从圆盘 $D$ 到 $\R^3$ 的 Lipschitz 分段 $C^1$ 映射且有一条
折叠曲线, 则存在分段 $C^1$ Lipschitz 变分向量场 $W$, 它是面积递减的.
特别地, $f$ 不可能是面积最小化的.
\end{lemma}

\begin{proof}
设 $f$ 面积最小. 则在它是浸入之处平均曲率为零. 折叠曲线的定义使我们
可以沿折痕构造一个面积递减的变分向量场 $W$.

由于 $f$ 在 $D$ 上 Lipschitz, 且 $f$ 在 $\{x_1\geq0\}$ 或
$\{x_1\leq0\}$ 上的限制是直到边界的 $C^1$ 浸入, $f_{x_2}$ 处处有定义.
另一方面,
$\lim_{x_1\to0_+}f_{x_1}(x_1,x_2)$ 与
$\lim_{x_1\to0_-}-f_{x_1}(x_1,x_2)$ 之间的夹角小于 $\pi$. 因此可以选取
紧支的 Lipschitz 向量场 $W$, 它是分段 $C^1$ 的, 使得一方面
\begin{equation}
\lim_{x_1\to0_+}\bigl\langle W(f(x_1,x_2)),f_{x_1}(x_1,x_2)\bigr\rangle>0,
\label{eq:6.128}
\end{equation}
另一方面
\begin{equation}
\lim_{x_1\to0_-}\bigl\langle W(f(x_1,x_2)),-f_{x_1}(x_1,x_2)\bigr\rangle>0,
\label{eq:6.129}
\end{equation}
最后
\begin{equation}
\bigl\langle W(f(0,x_2)),f_{x_2}\bigr\rangle=0.
\label{eq:6.130}
\end{equation}
于是立即看出, 若取 $W$ 为变分向量场, 则面积的一阶变分为负.
\end{proof}

引理~\ref{lem:6.35} 降低一维的类比是: 给定一条分段 $C^1$ 曲线
$\gamma$, 其切向 $\gamma'$ 在内部有间断, 则存在保持端点不动的单参数变形
使长度减小. 由此立即推出, 完备 Riemann 流形中两点间极小化测地线不可能
有这样的 ``折痕''.

我们先讨论在 $u(D)$ 自交为横截的特殊情形下如何用折叠曲线的概念来证明
嵌入性结果. 一般情形的证明将把这些想法与一个扰动论证结合起来.

我们需要的下一个结果是一个关于圆盘的自横截映射的纯拓扑结果. 设
$v:\bar{D}\to\R^3$ 是处于一般位置的 $C^1$ 真浸入 (不必极小), 且 $v$ 在
$\partial D$ 上的限制是嵌入.

我们仍可按显然的方式定义双重点集 (即有多于一个原像的点). 由于 $v$ 处于
一般位置, $v$ 的像沿有限多条浸入曲线之并 $\mathcal{D}$ 自交. 以
$U_1,\ldots,U_n$ 记 $D\setminus v^{-1}(\mathcal{D})$ 的连通分支, 于是
$v$ 在每个 $U_j$ 上的限制是嵌入.

圆盘 $D$ 由各 $\overline{U_j}$ 的不交并, 连同对边界 $\partial U_j$ 的一系列
等同而给出. 这些等同与 $v$ 相容, 其含义是: 在任意两个被等同的点处 $v$
取相同的值.

下面命题~\ref{prop:6.37} 的要点在于: 还存在另一种等同各 $\partial U_j$
的方式, 以同样的方式构造出 $D$, 它与 $v$ 相容, 并且使得 $v$ 现在是一个
分段 $C^1$ 嵌入, 只是在 $\mathcal{D}$ 处不再是 $C^1$ 的. 这样做之后,
新映射沿 $\mathcal{D}$ 有自交, 但不再自我穿越. 通过这一操作, 我们沿
$\mathcal{D}$ 引入了折叠曲线. 降低一维的下例说明了这一点:

\begin{example}
\label{ex:6.36}
设 $v:[-\pi,\pi]\to\R^2$ 由
\begin{equation}
v(t)=(2t/\pi-\sin t,\cos t)
\label{eq:6.131}
\end{equation}
给出. 则双重点集为 $(0,0)$, 其原像为
$\{-\frac{\pi}{2},\frac{\pi}{2}\}$, 并且
$U_1=(-\pi,-\frac{\pi}{2})$、$U_2=(-\frac{\pi}{2},\frac{\pi}{2})$、
$U_3=(\frac{\pi}{2},\pi)$. 起初, $[-\pi,\pi]$ 由取各 $U_i$ 的 (相对)
闭包之并, 并把 $-\frac{\pi}{2}\in\overline{U_1}$ 与
$-\frac{\pi}{2}\in\overline{U_2}$ 等同、把
$\frac{\pi}{2}\in\overline{U_2}$ 与
$\frac{\pi}{2}\in\overline{U_3}$ 等同而得到.

另一方面, 我们也可以通过把 $-\frac{\pi}{2}\in\overline{U_1}$ 与
$\frac{\pi}{2}\in\overline{U_2}$ 等同、把
$-\frac{\pi}{2}\in\overline{U_2}$ 与
$\frac{\pi}{2}\in\overline{U_3}$ 等同来构造区间 $[-\pi,\pi]$. 按这种
等同, $v$ 仍然连续, 因为
$v(-\frac{\pi}{2})=v(\frac{\pi}{2})$ (这就是我们所说的等同与 $v$ 相容).
按构造, $v$ 的像不变.

显然, $v$ 在 $-\frac{\pi}{2}$ 与 $\frac{\pi}{2}$ 处不再是 $C^1$ 的.
然而 $v$ 的像不再自我穿越, 尽管它在某个奇点处仍有一个自交 (这就是我们
所说的 $v$ 是一个分段 $C^1$ 嵌入). 现在我们可以通过在 $(0,0)$ 附近作
任意小的 $C^0$ 形变把 $v$ 变形成嵌入. 这一操作称为 ``磨圆'', 它将减小
$v$ 的像的长度.
\end{example}

我们将需要下述拓扑命题. 其证明见 Freedman \cite{Fr}, 用到了拓扑中的
塔式构造. 这一构造可追溯到 Papakyriakopolous \cite{Pa}.

\begin{proposition}
\label{prop:6.37}
设 $v:\bar{D}\to\R^3$ 是处于一般位置的 $C^1$ 真浸入, 且 $v$ 在
$\partial D$ 上的限制是嵌入. 若 $\mathcal{D}$ 是如上所述的双重点集, 以
$U_1,\ldots,U_n$ 记 $D\setminus v^{-1}(\mathcal{D})$ 的连通分支, 则通过
构造适当的与 $v$ 相容的等同, 可以构成一个新圆盘
$D'=\bigcup_{j=1}^{n}\overline{U}_j$, 具有下述性质. 由
$w|_{U_i}=v|_{U_i}$ 给出的映射 $w:\overline{D'}\to\R^3$ 是 Lipschitz
连续的分段 $C^1$ 嵌入. 映射 $w$ 与 $v$ 有相同的像, 且在
$\partial D=\partial D'$ 上一致. 最后, $w$ 在 $C^0$ 意义下任意接近一个
嵌入, 并且 $\partial U_i\setminus\partial D$ 的每个连通分支都是折叠曲线.
\end{proposition}

在 $u(D)$ 自横截的情形下, 定理~\ref{thm:6.29} 的证明现在已经清楚了.
也就是说, 我们用命题~\ref{prop:6.37} 为一个具有相同面积与边界值的映射
产生一条折叠曲线. 引理~\ref{lem:6.35} 蕴涵这个新映射不是面积最小化的,
从而 $u$ 本来也不是. 这一矛盾蕴涵 $u$ 一开始就必是嵌入.

利用 Freedman \cite{Fr} 以及 Freedman、Hass 与 Scott \cite{FHS} 的想法,
可以把上述途径与一个扰动论证结合起来以处理一般情形.

\begin{lemma}
\label{lem:6.38}
对如上所述的 $u$ 与 $\mathcal{D}$, 存在 $\epsilon>0$ 与
$x_1\in\mathcal{D}$, 使得 $B_\epsilon(x_1)\cap\mathcal{D}$ 只有横截
自交.
\end{lemma}

\begin{proof}
这直接由非横截自交是孤立的推出.
\end{proof}

在继续之前, 我们回忆浸入极小曲面在自交点附近的局部描述. 这一结果是
关于椭圆方程零点集与临界集的解析结果的直接推论 (类似于\secref{sec:6.2}中的论证).

\begin{lemma}[Freedman, Hass 与 Scott, \cite{FHS}]
\label{lem:6.39}
设 $U\subset D$ 是 $0$ 的邻域, $u_1,u_2$ 是 $U$ 上的光滑函数, 并设
$v=u_1-u_2$ 满足 $v=c\operatorname{Re}z^d+q(x)$, 其中 $c\neq0$ 且 $q$
满足 \eqref{eq:6.39}. 设 $h:U\to V$ 是 $C^1$ 微分同胚, 其中
$V\subset D$ 是 $0$ 的邻域, 且对 $x\in U$ 有
$c\operatorname{Re}z^d(h(x))=v(x)$.

取 $r>0$ 使 $D_r\subset U$, 并设 $\psi$ 是 $U$ 上支集含于 $D_r$ 的
光滑函数, 满足当 $|x|\leq\frac{r}{2}$ 时 $\psi(x)=1$. 对 $t>0$ 置
$u_t(x)=u_1(x)+t\psi(x)$; 则存在 $t_0>0$, 使得当 $0<t\leq t_0$ 时
$u_t$ 与 $u_2$ 的图横截相交.
\end{lemma}

\begin{remark}
\label{rem:6.40}
注意: 由于在 $\partial D_r$ 之上相交情况不变, 新的相交曲线与原 (奇异)
曲线有相同的端点.
\end{remark}

\mtpara{嵌入性}
现在我们准备证明本节的主要结果.

\begin{proof}[定理~\ref{thm:6.29} 的证明]
到此为止, 我们已经证明 Plateau 问题的解 $u:D\to\Omega$ 是光滑真浸入,
且在边界附近是嵌入. 还需证明双重点集 $\mathcal{D}$ 为空. 我们假设
$\mathcal{D}\neq\emptyset$ 并由此导出矛盾.

由引理~\ref{lem:6.38}, 存在 $\epsilon_1>0$ 与
$x_1\in\mathcal{D}$, 使 $B_{\epsilon_1}(x_1)\cap\mathcal{D}$ 只有横截
自交. 由推论~\ref{cor:6.33}, 为使 $B_{\epsilon_1}(x_1)\cap u(D)$ 成为
嵌入, 只有有限多种可能的剪贴操作 (如命题~\ref{prop:6.37} 中那样). 由
引理~\ref{lem:6.35}, 这些剪贴操作中的每一个都至少使 $u(D)$ 的面积减少
某个 $\epsilon_2>0$ (通过沿所得折叠曲线作扰动).

引理~\ref{lem:6.30} 蕴涵存在某个 $\delta>0$ 使
$u(D_{1-\delta})\cap\mathcal{D}=\emptyset$. 因此, 局部描述
引理~\ref{lem:6.39} 使我们可以在
\begin{equation}
D_{1-\delta}\setminus u^{-1}\bigl(B_{\epsilon_1}(x_1)\bigr)
\label{eq:6.132}
\end{equation}
中作任意小的扰动, 以产生一个处于一般位置的新映射 $\tilde{u}$. 特别地,
可以假设
\begin{equation}
\Area\bigl(\tilde{u}(D)\bigr)<\Area\bigl(u(D)\bigr)+\frac{\epsilon_2}{2}.
\label{eq:6.133}
\end{equation}
由命题~\ref{prop:6.37}, 我们可以剪贴 $D$ 以得到新圆盘 $D'$, 使得这些
等同与 $\tilde{u}$ 相容. 分段光滑嵌入 $w:D'\to\Omega$ 与 $\tilde{u}$ 有
相同的像, 从而面积相同, 且这两个映射在边界上一致. 由于它们的像相同,
在 $u^{-1}(B_{\epsilon_1}(x_1))$ 中所作的等同必属于上面列举的那些.
特别地, 通过沿 $B_{\epsilon_1}(x_1)$ 中所得的某条折叠曲线作扰动, 我们
可以使 $w(D')$ 的面积至少减少 $\epsilon_2$. 结合 \eqref{eq:6.133},
这与 $\Area(u(D))$ 的极小性矛盾, 证明完成.
\end{proof}
